%% ════════════════════════════════════════════════════════════════════════
%% CRF_Complete_v17_16_Part_L.tex
%% Complete Edition v17.16 — Part L (ZC88–93)
%%
%% Two clusters:
%%   Cluster 1 (Force-Layer Bridge): ZC88 Log-Bridge, ZC89 Layer-Specificity
%%   Cluster 2 (Mind-Layer Structure): ZC90 Quantum-Pair, ZC91 Irrationality-
%%     of-Kind, ZC92 Barami Functional, ZC93 Gate-Ratio
%%
%% Governance: Upgrade Protocol v1.3 · CRF Style Guide v3.7 (UPG-30/31) ·
%%   Integration Execution Script v1.2 · Lexicon Lock v4 · Governance Manifest v1.
%%   Part L is the first v3.7-native Part (v3.7 applies from ZC88 onward;
%%   Parts A–K were built under v3.6/earlier).
%%
%% No-Loss: all six ZC bodies embedded verbatim EXCEPT label/ref namespacing
%%   (\label{X}->\label{X-zcNN}, refs rewritten to match) to remove cross-paper
%%   collisions. Macro set-diff per Script v1.2 STEP A2; new macros in
%%   CRF_v17_16_macro_patch.tex (consolidated, single-file, no \input chain).
%%
%% Discovery origins: all six papers carry %% Discovery origin: fields.
%% ════════════════════════════════════════════════════════════════════════

\documentclass[12pt, a4paper]{article}

%% ── PACKAGES ────────────────────────────────────────────────────────────
\usepackage{amsmath, amssymb, amsthm}
\usepackage{mathtools}
\usepackage{geometry}
\usepackage{xcolor}
\usepackage{parskip}
\usepackage[protrusion=false,expansion=false]{microtype}
\usepackage{hyperref}
\usepackage{tcolorbox}
\usepackage{booktabs}
\usepackage{longtable}
\usepackage{array}
\usepackage[T1]{fontenc}
\usepackage{graphicx}
\usepackage{enumitem}
\usepackage{needspace}
\usepackage{tocloft}

%% Unicode-engine font support; morewrites prevents \write exhaustion.
\usepackage{iftex}
\iftutex
  \usepackage{morewrites}
  \usepackage{fontspec}
  \setmainfont{Latin Modern Roman}
  \usepackage{polyglossia}
  \setdefaultlanguage{english}
  \setotherlanguage{thai}
  \newfontfamily\thaifont[Scale=0.95]{Loma}
\fi

\geometry{margin=2.5cm}

%% TOC spacing
\setlength{\cftbeforesecskip}{0pt}
\setlength{\cftbeforesubsecskip}{0pt}
\setlength{\cftbeforepartskip}{6pt}
\setlength{\cftsubsecindent}{2.5em}
\setlength{\cftsubsubsecindent}{4.5em}
\renewcommand{\cftsecfont}{\normalfont}
\renewcommand{\cftsubsecfont}{\small}
\renewcommand{\cftsecpagefont}{\normalfont}
\renewcommand{\cftsubsecpagefont}{\small}

%% ── COLOR PALETTE (Style Guide v3.5) ────────────────────────────────────
\definecolor{deepblue}{RGB}{10,30,80}
\definecolor{crfgreen}{RGB}{0,100,60}
\definecolor{softgreen}{RGB}{180,230,210}
\definecolor{physgold}{RGB}{140,100,0}
\definecolor{softgold}{RGB}{255,240,180}
\definecolor{loopred}{RGB}{160,30,30}
\definecolor{softred}{RGB}{255,210,210}
\definecolor{mutedgray}{RGB}{90,90,90}
\definecolor{midblue}{RGB}{30,80,160}
\definecolor{softblue}{RGB}{180,210,240}
\definecolor{layerblue}{RGB}{20,60,120}
\definecolor{genealogygold}{RGB}{120,80,0}
\definecolor{softgenealogy}{RGB}{255,248,220}
\definecolor{conmapbg}{RGB}{240,245,255}
\definecolor{conmapframe}{RGB}{50,80,140}
\definecolor{inheritbg}{RGB}{225,240,255}
\definecolor{inheritframe}{RGB}{20,60,160}
\definecolor{correctbg}{RGB}{255,235,220}
\definecolor{correctframe}{RGB}{180,90,30}
\definecolor{giftgold}{RGB}{180,140,0}
\definecolor{softgift}{RGB}{255,252,200}
\definecolor{nobelblue}{RGB}{10,40,100}
\definecolor{softnobel}{RGB}{200,220,255}

\hypersetup{colorlinks=true, linkcolor=deepblue, urlcolor=midblue,
            citecolor=deepblue, bookmarks=false}

%% ── BOX STYLES (Style Guide v3.5 — all 10 types) ────────────────────────
\tcbuselibrary{skins, breakable}

\newtcolorbox{derivedbox}[1][]{colback=softgreen!50,
  colframe=crfgreen!60, fonttitle=\bfseries\color{crfgreen!20!black},
  title={Derived}, breakable, #1}
\newtcolorbox{formalbox}[1][]{colback=softblue!40,
  colframe=midblue!60, fonttitle=\bfseries\color{deepblue},
  title={Formal}, breakable, #1}
\newtcolorbox{openqbox}[1][]{colback=softred!30,
  colframe=loopred!50, fonttitle=\bfseries\color{loopred!20!black},
  title={Open}, breakable, #1}
\newtcolorbox{keybox}[1][]{colback=softgold!50,
  colframe=physgold!60, fonttitle=\bfseries\color{physgold!20!black},
  breakable, #1}
\newtcolorbox{layerbox}[1][]{
  colback=layerblue!8, colframe=layerblue!50,
  fonttitle=\bfseries\color{layerblue},
  title={R-Layer Note}, breakable, #1}
\newtcolorbox{genealogybox}[1][]{
  colback=softgenealogy!60, colframe=genealogygold!70,
  fonttitle=\bfseries\color{genealogygold!30!black},
  title={Genealogy Map}, breakable, #1}
\newtcolorbox{conmapbox}[1][]{
  colback=conmapbg, colframe=conmapframe,
  fonttitle=\bfseries\color{white},
  breakable, #1}
\newtcolorbox{inheritbox}[1][]{
  colback=inheritbg, colframe=inheritframe,
  fonttitle=\bfseries\color{white},
  title={Backward Inheritance}, breakable, #1}
\newtcolorbox{correctionbox}[1][]{
  colback=correctbg, colframe=correctframe,
  fonttitle=\bfseries\color{white},
  title={Forward Correction / Cross-Part PM}, breakable, #1}
\newtcolorbox{contentsbox}[1][]{colback=softblue!20,
  colframe=midblue!50, fonttitle=\bfseries\color{deepblue},
  breakable, #1}
%% Style Guide v3.5: giftbox (ZC65+) and nobelbox (ZC70+)
\newtcolorbox{giftbox}[1][]{
  colback=softgift!70, colframe=giftgold!80,
  fonttitle=\bfseries\color{giftgold!20!black},
  title={Gift}, breakable, #1}
\newtcolorbox{nobelbox}[1][]{
  colback=softnobel!60, colframe=nobelblue!70,
  fonttitle=\bfseries\color{nobelblue},
  title={Nobel-Table Claim}, breakable, #1}

%% ── Theorem environments ────────────────────────────────────────────────
\newtheorem{theorem}{Theorem}[section]
\newtheorem{lemma}{Lemma}[section]
\newtheorem{corollary}{Corollary}[section]
\newtheorem{finding}{Finding}[section]
\newtheorem{proposition}{Proposition}[section]
\newtheorem{definition}{Definition}[section]
\newtheorem{principle}{Principle}[section]
\newtheorem{claim}{Claim}[section]
\newtheorem{conjecture}{Conjecture}[section]
\newtheorem*{remark}{Remark}

%% volumebox (Part overview summary)
\newtcolorbox{volumebox}[1][]{colback=softblue!15,
  colframe=deepblue!50, fonttitle=\bfseries\color{deepblue},
  breakable, #1}

%% ── Title ───────────────────────────────────────────────────────────────
\title{%
  \textbf{\color{deepblue}CRF Complete Edition v17.16 --- Part L}\\[0.4em]
  \large ZC88--93: The Log-Bridge and the Mind-Layer\\[0.2em]
  \normalsize\color{mutedgray}
  Force-Layer Bridge (log-unification, layer-specificity) $\cdot$
  Mind-Layer Structure (quantum-pair, irrationality-of-kind,
  Barami functional, gate-ratio)\\[0.2em]
  \textbf{\color{crfgreen}Two mirrors shown one law in log-space;
  the mind-layer given its own discrete--continuous pair}
}
\author{Ougit Jarittum\\
\small Independent Researcher, Thailand\quad
ORCID: 0009-0009-4451-6811\\
\small Zenodo: \href{https://doi.org/10.5281/zenodo.18977059}%
{10.5281/zenodo.18977059}\quad (CC-BY-NC 4.0)
}
\date{June 2026 --- Complete Edition v17.16 (Part L of 12)}

%% ── PART COUNTER: continue from Part K (ends at LVI = 56) ──────────────
\setcounter{part}{56}

%% ── MACRO PATCH v17.16 (consolidated; replaces v17_14 + v17_15 chain) ───
\input{CRF_v17_16_macro_patch.tex}

%% ── Part L local notation overrides (No-Loss, edition-level) ──────────────
%% ZC88–93 (Style Guide v3.7) display the ε-family bare, matching ZC83–87.
%% The consolidated chain tags some ε macros with zero-subscript / sector
%% forms; override to bare here. Chain + sources unmodified. Mirrors the
%% Part K PM-V17_15-EPSFAMILY override exactly.
\renewcommand{\epsEM}{\varepsilon_{\rm EM}}     % chain: ε_0^{EM}
\renewcommand{\epscan}{\varepsilon_{\rm can}}   % chain: ε_0^{can}
\renewcommand{\epsSC}{\varepsilon_{\rm SC}}     % chain: ε_{can}^{SC}
\renewcommand{\epsS}{\varepsilon_{\rm strong}}  % chain first-wins: ε_{0,Spont}

%% ════════════════════════════════════════════════════════════════════════
\begin{document}
\sloppy

%% Governance: Upgrade Protocol v1.3, CRF Style Guide v3.7 (UPG-30/31),
%%             Integration Execution Script v1.2, Lexicon Lock v4

\maketitle
\tableofcontents
\newpage

%% ════ Reader's On-Ramp (Protocol v1.3 §U2; Style Guide v3.7 §31.2) ═══════
%% Primary audience tier: 1   |   On-ramp for tier below (Tier 0): yes
\begin{keybox}[title={If you are just arriving --- Part L in plain language}]
\sloppy
CRF tries to derive the structure of physics from a single starting idea
called Distinction ($\delta$). Two earlier threads each found a ``mirror'':
the force sectors' instability floors sit in a multiplicative pattern (each
floor a fixed ratio from the next), and a separate ladder of rates sits in
an additive pattern (each rung a fixed step from the next). \textbf{Cluster
1} (ZC88--89) shows these are not two facts but \emph{one}: take the
logarithm and the multiplicative mirror becomes the additive one --- same
law, two disguises --- then pins down exactly which physical layer each step
belongs to. \textbf{Cluster 2} (ZC90--93) carries the same discrete-versus-
continuous idea into the ``mind layer'': it pairs a smooth step with a
sudden one, shows the two kinds cannot be converted into each other, and
builds from them a capacity functional (Barami) and the gate that controls
which kind of step a choice can make.

\textbf{Minimum vocabulary:}
\textit{$\delta$ (Distinction)} --- the one starting idea;\;
\textit{floor / $\varepsilon$} --- an instability threshold, one per force
sector;\;
\textit{log-bridge} --- taking $\ln$ turns a multiplicative pattern into an
additive one;\;
\textit{visesa / nibbedha} --- the smooth (continuous) and sudden
(discontinuous) mind-steps;\;
\textit{Barami} --- a capacity built up from accumulated forward choices;\;
\textit{gap / GAP-\dots} --- an open question the chain has not yet closed.
Full symbol list: Master Notation Key (front matter).

These results are pieces of a map under construction, not a claim of final
truth: each carries the conditions under which it would fail (see front
matter, ``Map, not reality'').
\end{keybox}

\newpage

%%=============================================================
\part{Force-Layer Bridge: Log-Unification and Layer-Specificity (v17.16)}
\label{part:LVII}
%%=============================================================

% [BRIDGE-PROSE: integration-added, Style Guide v3.7 §31.6 / UPG-30]
\noindent
By the end of Part K the force sectors had two separate descriptions that
nobody had connected. The floor lattice put the instability thresholds in a
\emph{multiplicative} order --- each floor a fixed ratio above the last, with
one floor sitting exactly at the geometric mean of two others. The
$\Gamma$-ladder put a family of rates in an \emph{additive} order --- each
rung a fixed amount above the last. The stakes here are not cosmetic: if
these two orders were genuinely independent, CRF would carry two unrelated
structural laws where it claims to have one origin. ZC88 asks whether the
logarithm collapses the difference, and what survives if it does.
\section{Background: Two Mirrors, One Suspected Ladder}
\label{sec:background-zc88}
%%=============================================================

%% ── READER'S ON-RAMP (v3.6 §31.2 / UPG-25) ────────────────
\begin{keybox}[title={If you are just arriving --- ZC88 in plain language}]
\sloppy
CRF derives the constants of physics from one primitive, Distinction
($\delta$). Two of its results are short lists of numbers with a built-in
symmetry. The \emph{floor list} $(\epsEM,\epsuniv,\epsW)$ describes the
smallest instability each force sector can carry; its middle value is the
\emph{geometric} mean of the outer two. The \emph{rate list}
$(\GammaS,\GammaW,\GammaEM)$ describes how fast each sector runs; its
middle value is the \emph{arithmetic} mean of the outer two. This paper
shows that these are not two coincidences but one: take the logarithm of
the floor list and it becomes an evenly spaced list with the same shape as
the rate list. ``Geometric mean'' and ``arithmetic mean'' are the same
statement seen through a logarithm.

\textbf{Minimum vocabulary:}
\textit{$\epsilon$ floor} --- minimum sector instability;\;
\textit{$\Gamma$ rung} --- sector running rate;\;
\textit{AP} --- arithmetic progression (evenly spaced);\;
\textit{GM/AM} --- geometric / arithmetic mean.
CRF is a map under construction, not a claim of final truth.
\end{keybox}

%% ── NUMERICAL SUMMARY ──────────────────────────────────────
\begin{keybox}[title={Numerical Summary --- ZC88 (T-canonical)}]
{\small\renewcommand{\arraystretch}{1.25}
\begin{center}
\begin{tabular}{@{}p{5.4cm}p{2.6cm}p{4.4cm}@{}}
\toprule
\textbf{Quantity} & \textbf{Value} & \textbf{Source / Status} \\
\midrule
$r=\epsW/\epsEM$              & $1.028085$   & THM-ZC67-01 \\
$\tfrac14\ln r$ (floor log-step) & $0.00692457$ & LEM-ZC88-01 (NEW) \\
$\DeltaGam=\ln 2/\ln 15$      & $0.25595802$ & FIND-ZC73-01 \\
$\tfrac12(\ln\epsEM+\ln\epsW)$ & $-4.8797734$ & $=\ln\epsuniv$ (COR-ZC88-01) \\
$\ln\epsuniv$                 & $-4.8797734$ & exact (diff $<10^{-15}$) \\
$\DeltaGam/(\tfrac14\ln r)$   & $36.9638$    & \textbf{NEW: FIND-ZC88-01 (not clean)} \\
\bottomrule
\end{tabular}
\end{center}}
\end{keybox}

\subsection{What was already known before ZC88}

\begin{keybox}[title={Pre-ZC88 Status of the Floor / $\Gamma$ Structures}]
\textbf{Known:}
\begin{itemize}[leftmargin=2em, noitemsep]
  \item LEM-ZC83-01 (Part~K): the floor lattice is
    $\epsEM=\epsuniv\,r^{-1/4}$, $\epsW=\epsuniv\,r^{+1/4}$, with $r>1$.
  \item FIND-ZC83-01 (Part~K): $\sqrt{\epsEM\,\epsW}=\epsuniv$ exactly
    --- the bridge floor is the geometric mean of the sector floors.
  \item THM-ZC73-01 / FIND-ZC81-01 (Parts~I/J): the $\Gamma$-ladder
    $\Gamma_k=\Gcoup/\ds-1+k\,\DeltaGam$ ($k=0,1,2,3$) is an arithmetic
    progression of step $\DeltaGam=\ln 2/\ln 15$, with mirror axis
    $\GammaW$ ($k=2$): $\GammaEM+\GammaS=2\,\GammaW$.
\end{itemize}

\textbf{Open (before ZC88):}
\begin{itemize}[leftmargin=2em, noitemsep]
  \item Whether the multiplicative floor mirror and the additive $\Gamma$
    mirror are independent facts or a single structure.
  \item Whether the two ladder step-sizes are related.
\end{itemize}

ZC88 answers the first (they are one structure) and bounds the second
(the steps are \emph{not} simply related; GAP-ZC88-01).
\end{keybox}

\begin{keybox}[title={What ZC88 Does NOT Do}]
\begin{itemize}[leftmargin=2em, noitemsep]
  \item Does not change any $\epsilon$ floor value or any $\Gamma$ rung
    value (all inherited verbatim).
  \item Does not claim $\DeltaGam$ and $\tfrac14\ln r$ are proportional
    --- the data refute it (FIND-ZC88-01).
  \item Does not revisit the FIND-ZC81-01 or FIND-ZC83-01 proofs
    (used as pillars, stand unchanged).
  \item Does not touch $\epsSC$ ontology (GAP-ZC83-01, closed in Part~K).
\end{itemize}
\end{keybox}

\begin{keybox}[title={Companion Papers --- ZC88}]
{\small\renewcommand{\arraystretch}{1.25}
\begin{center}
\begin{tabular}{@{}p{5.2cm}p{7.3cm}@{}}
\toprule
\textbf{Paper} & \textbf{What it provides to ZC88} \\
\midrule
\texttt{Part~K (ZC83)}
  & LEM-ZC83-01 floor formula $r^{\pm1/4}$; FIND-ZC83-01 GM identity. \\
\texttt{Part~J (ZC81/82)}
  & FIND-ZC81-01 mirror $\GammaEM+\GammaS=2\GammaW$; AP step $\DeltaGam$. \\
\texttt{Part~I (ZC73)}
  & THM-ZC73-01 $\Gamma$-ladder; FIND-ZC73-01 $\DeltaGam=\ln2/\ln15$. \\
\texttt{CRF\_Complete\_Index}
  & Parts~J/K quick-reference: floor + $\Gamma$ values. \\
\bottomrule
\end{tabular}
\end{center}}
\end{keybox}

\begin{derivedbox}[title={Discovery Note}]
This paper began in the Meta-Phase $\Omega$ corpus survey (GIFT-$\Omega$-01).
Seen separately, the floor mirror (multiplicative) and the $\Gamma$ mirror
(additive) look like two unrelated gifts. Laid side by side and fed through
a logarithm, the floor lattice returned an arithmetic progression with axis
$\ln\epsuniv$ to machine precision --- visibly the same shape as the
$\Gamma$-ladder. The proof direction became clear at once: the geometric
mirror \emph{is} the arithmetic mirror, read on a log scale. The honest
caveat surfaced in the same probe: the two step-sizes do not match, so the
bridge is structural, not metric.
\end{derivedbox}

\subsection{Pre-Work Audit (PWA-1 through PWA-10)}

\begin{formalbox}[title={Pre-Work Audit Record --- ZC88}]
Scanned before writing (Upgrade Protocol v1.3, PWA-1 through PWA-11):
\begin{itemize}[leftmargin=2em, noitemsep]
  \item LEM-ZC83-01 / FIND-ZC83-01 (Part~K) --- \textbf{KEY PILLAR}
    (floor formula + GM identity, used verbatim).
  \item FIND-ZC81-01 / THM-ZC73-01 (Parts~J/I) --- \textbf{KEY PILLAR}
    ($\Gamma$ mirror + ladder step, used verbatim).
  \item Index Quick-Reference Pack $\varepsilon$-floor notice scanned:
    the three floor values are one identity on three input pairs; ZC88
    introduces \emph{no} new value and does not ``reconcile'' them.
\end{itemize}

\textbf{PWA-6} (macro convention): six new $\Gamma$/mean macros added by
\texttt{\textbackslash providecommand}; no conflict with Preamble~v3.1.\\
\textbf{PWA-7} (sector-floor): not triggered (no floor reassigned).\\
\textbf{PWA-8} (layer self-consistency): all objects at the $R_0$ layer;
$\epsuniv$ is the bridge floor, $\epscan$ not used.\\
\textbf{PWA-9} (spectral convention): not triggered.\\
\textbf{PWA-10} (tautology audit): the log-bridge is \emph{not} a
restatement of one definition --- it relates two independently derived
sub-chains (Born-chain floors vs.\ K14 fixed point vs.\ $\Gamma$-ladder).
See PM-ZC88-01 for the refuted proportionality probe.\\
\textbf{No duplicate atomic units found.}
\end{formalbox}


% [BRIDGE-PROSE: integration-added, Extension Freedom narrative, v3.6 §31.3]
CRF had accumulated two gifts sitting in adjacent shelves without anyone
placing them next to each other. The floor lattice --- four instability
thresholds, one per force sector --- obeys a geometric-mean symmetry: the
middle value is the exact square root of the outer two. The $\Gamma$-ladder
--- the four running rates of the same sectors --- obeys an arithmetic-mean
symmetry: the middle rung is the exact average of the outer two. Two
different kinds of mean, two different lists of numbers, two different papers.
The question nobody had asked was whether they are two symmetries or one.

The answer turns on what happens when you take the logarithm of the floor
list. Logarithm converts multiplication into addition; geometric mean into
arithmetic mean. If the floor symmetry and the rate symmetry were truly
separate, the log-floor list would have no particular relationship to the
rate list. What the numerical probe found instead was that the log-floor
list is an evenly spaced sequence --- an arithmetic progression --- with
exactly the same three-term mirror structure as the rate ladder. The two
shelves hold the same object. The sections below prove it in two lines.

\newpage

%%=============================================================
\section{The Floor Lattice in Log-Space\quad (LEM-ZC88-01)}
\label{sec:lem-zc88}
%%=============================================================

\needspace{8\baselineskip}

\begin{formalbox}[title={LEM-ZC88-01: The Log Floor Lattice is an
  Arithmetic Progression ($R_0$, Exact)}]

\textbf{Statement.}
The logarithms of the three floor values form an arithmetic progression
with mirror axis $\ln\epsuniv$:
\begin{equation}
  \boxed{
    \ln\epsEM,\;\ln\epsuniv,\;\ln\epsW
    \quad\text{is an AP of step }s_{\rm floor}=\tfrac14\ln r,
    \quad r=\frac{\epsW}{\epsEM}.
  }
  \label{eq:lem-main-zc88}
\end{equation}

\textbf{Proof.}

\textit{Step 1 (substitute the floor formula).}
From LEM-ZC83-01, $\epsEM=\epsuniv\,r^{-1/4}$ and
$\epsW=\epsuniv\,r^{+1/4}$. Taking logs,
\[
  \ln\epsEM=\ln\epsuniv-\tfrac14\ln r,\qquad
  \ln\epsW =\ln\epsuniv+\tfrac14\ln r.
\]

\textit{Step 2 (read off the AP).}
The three terms are $\ln\epsuniv-s_{\rm floor}$, $\ln\epsuniv$,
$\ln\epsuniv+s_{\rm floor}$ with $s_{\rm floor}=\tfrac14\ln r$. Consecutive
differences are equal, so the sequence is an arithmetic progression; the
middle term $\ln\epsuniv$ is its axis of symmetry.
\hfill$\square$

\textbf{Numerical verification (T-canonical):}
\begin{center}
{\small\renewcommand{\arraystretch}{1.25}
\begin{tabular}{@{}lrl@{}}
\toprule
Quantity & Value & Source \\
\midrule
$r=\epsW/\epsEM$            & $1.028085$    & THM-ZC67-01 \\
$s_{\rm floor}=\tfrac14\ln r$ & $0.00692457$ & this lemma \\
$\ln\epsW-\ln\epsuniv$     & $0.00692457$  & $=+s_{\rm floor}$ \\
$\ln\epsuniv-\ln\epsEM$    & $0.00692457$  & $=+s_{\rm floor}$ \\
\bottomrule
\end{tabular}}
\end{center}

\textbf{Status:} Exact Lemma.
\physmeaning{the floor lattice is multiplicative on the surface, but its
  ``true'' coordinate is the logarithm: in that coordinate it is a perfectly
  even three-rung ladder, exactly like the $\Gamma$-ladder.}
\end{formalbox}


% [BRIDGE-PROSE: integration-added, Extension Freedom narrative, v3.6 §31.3]
The lemma settles what the floor list looks like in log-coordinates: an
evenly spaced three-rung ladder with the bridge floor at the centre. The
$\Gamma$-ladder already has that shape in linear coordinates. Saying the two
are ``co-structural'' is therefore a precise claim, not an analogy --- both
are instances of the same abstract object, a mirror-symmetric arithmetic
progression, one read on a log scale and one on a linear scale. The theorem
below states this formally and maps the bridge between them.

%%=============================================================
\section{The Co-Structure Theorem\quad (THM-ZC88-01)}
\label{sec:thm-zc88}
%%=============================================================

\needspace{10\baselineskip}

\begin{formalbox}[title={THM-ZC88-01: Floor Lattice and $\Gamma$-Ladder
  are Co-Structural ($R_0$, Exact)}]

\textbf{Statement.}
Both structures are three-term mirror-symmetric arithmetic progressions
about a middle term. Writing each as
$(\,c-s,\;c,\;c+s\,)$:
\begin{equation}
  \boxed{
  \begin{aligned}
    \text{floor (log):}\quad
      &\big(\ln\epsEM,\;\ln\epsuniv,\;\ln\epsW\big)
       =\big(c_{\rm f}-s_{\rm f},\,c_{\rm f},\,c_{\rm f}+s_{\rm f}\big),\\
    \text{$\Gamma$-ladder:}\quad
      &\big(\GammaS,\;\GammaW,\;\GammaEM\big)
       =\big(c_\Gamma-s_\Gamma,\,c_\Gamma,\,c_\Gamma+s_\Gamma\big),
  \end{aligned}}
  \label{eq:thm-main-zc88}
\end{equation}
with $c_{\rm f}=\ln\epsuniv$, $s_{\rm f}=\tfrac14\ln r$ (LEM-ZC88-01) and
$c_\Gamma=\GammaW$, $s_\Gamma=\DeltaGam$ (THM-ZC73-01/FIND-ZC81-01). The
log-map is the bridge between the multiplicative face (floors) and the
additive face ($\Gamma$) of one ladder.

\textbf{Derivation.}

\textit{Step 1 (floor side).}
LEM-ZC88-01 gives the log-floor AP with axis $c_{\rm f}=\ln\epsuniv$.

\textit{Step 2 ($\Gamma$ side).}
THM-ZC73-01 gives $\Gamma_k=\Gcoup/\ds-1+k\,\DeltaGam$. Restricting to the
matter sectors $k=1,2,3$ (Strong, Weak, EM) yields
$(\GammaS,\GammaW,\GammaEM)=(c_\Gamma-\DeltaGam,\,c_\Gamma,\,c_\Gamma+\DeltaGam)$
with $c_\Gamma=\GammaW$. This is the same $(c-s,c,c+s)$ shape.

\textit{Step 3 (the bridge map).}
Multiplicative-on-floors becomes additive-on-logs:
$\epsilon\mapsto\ln\epsilon$ turns the geometric ratio $r^{\pm1/4}$ into the
additive offset $\pm\tfrac14\ln r$. The two ladders therefore live on
opposite sides of the exponential/logarithm pair; they are the same abstract
object (a mirror-symmetric 3-term AP) presented on two scales.

\textbf{Order note (No-Loss).}
The floor list is written low$\to$high ($\epsEM<\epsuniv<\epsW$, i.e.\
$k=3,?,?$ \emph{not} aligned to $\Gamma$'s $k$); the $\Gamma$ list is
written low$\to$high in $k$ ($\GammaS<\GammaW<\GammaEM$). The co-structure
claim is about \emph{shape} (mirror-symmetric AP), not about a sector-by-sector
index identification --- ZC81 already records $\epsS=\epsEM$ (spatial-type),
so the floor and rate sector-labels do not map one-to-one. PWA-7 honoured.
\hfill$\square$

\textbf{Status:} Exact Theorem (structural).
\physmeaning{the corpus has been carrying one ladder written twice: once as
  multiplicative floors, once as additive rates. Exponentiation and
  logarithm are the only difference between the two faces.}
\end{formalbox}

%%=============================================================
\section{Corollary: The Two Mirrors Are One Law\quad (COR-ZC88-01)}
\label{sec:cor-zc88}
%%=============================================================

% [BRIDGE-PROSE: integration-added, Extension Freedom narrative, v3.6 §31.3]
The theorem has established that two ladders share a shape. The corollary
pushes further and asks: not just a similar shape, but literally the same
equation? The answer is yes, and the proof fits on one line. Every step of
the argument is already in the literature --- the geometric-mean identity in
Part~K, the arithmetic-mean identity in Part~J, the logarithm as an
algebraic tautology. What was missing was the observation that these three
things are not separate facts but a single statement, and that making the
observation required no new machinery at all. The corollary below states
this collapse formally.

\needspace{8\baselineskip}

\begin{formalbox}[title={COR-ZC88-01: FIND-ZC83-01 $\equiv$ FIND-ZC81-01
  under the log-map ($R_0$, Exact)}]

\textbf{Statement.}
The geometric-mean floor mirror and the arithmetic-mean $\Gamma$ mirror are
the same identity, because the arithmetic mean of logarithms equals the
logarithm of the geometric mean:
\begin{equation}
  \boxed{\;\tfrac12\big(\ln\epsEM+\ln\epsW\big)=\ln\epsuniv\;}
  \label{eq:cor-main-zc88}
\end{equation}
This single line wears three faces, one per structure:
\begin{align*}
  \text{floor (GM, FIND-ZC83-01):}\quad
    &\sqrt{\epsEM\,\epsW}=\epsuniv,\\
  \text{its log (AM of logs):}\quad
    &\tfrac12\big(\ln\epsEM+\ln\epsW\big)=\ln\epsuniv,\\
  \text{$\Gamma$ (AM, FIND-ZC81-01):}\quad
    &\tfrac12\big(\GammaEM+\GammaS\big)=\GammaW.
\end{align*}

\textbf{Resolution.}
Apply $\ln(\cdot)$ to FIND-ZC83-01:
\[
  \ln\sqrt{\epsEM\,\epsW}
  =\tfrac12(\ln\epsEM+\ln\epsW)
  =\ln\epsuniv.
\]
The left form says ``the middle floor is the geometric mean of the outer
two''; the middle form says ``the axis of the log-AP is the arithmetic mean
of the outer two log-terms.'' These are literally the same equation.
FIND-ZC81-01 states exactly ``the middle rung is the arithmetic mean of the
outer two rungs'' for the $\Gamma$-ladder. By THM-ZC88-01 both ladders are
mirror-symmetric APs, so ``middle = mean'' is one structural law instantiated
on each face.

\textbf{Axiom chain.}
\[
  \begin{aligned}
  \delta \;\to\;& \mathbb{Z}_{15}\ \text{sector weights}
  \;\to\;\Big\{\,r^{\pm1/4}\ \text{(floors)},\ \ \pm\DeltaGam\ (\Gamma)\,\Big\}\\
  \;\to\;& \text{mirror-symmetric AP}
  \;\to\;\text{middle}=\text{mean}.
  \end{aligned}
\]

\textbf{Numerical verification (T-canonical):}
\begin{center}
{\small\renewcommand{\arraystretch}{1.25}
\begin{tabular}{@{}lrl@{}}
\toprule
Quantity & Value & \\
\midrule
$\tfrac12(\ln\epsEM+\ln\epsW)$ & $-4.87977336$ & AM of logs \\
$\ln\epsuniv$                  & $-4.87977336$ & axis \\
difference                     & $<10^{-15}$   & machine precision \\
$\tfrac12(\GammaEM+\GammaS)$   & $-0.24229356$ & AM of rungs \\
$\GammaW$                      & $-0.24229356$ & axis \\
\bottomrule
\end{tabular}}
\end{center}

\textbf{Status:} Exact Corollary (structural unification; no value changed).
\end{formalbox}


% [BRIDGE-PROSE: integration-added, Extension Freedom narrative, v3.6 §31.3]
The corollary collapses the apparent distance between the two mirror laws to
zero: they are literally the same equation, wearing different notation.
That is the satisfying half of this paper. The unsatisfying half --- and the
honest half --- is what happens next. Once you know two ladders share a
shape, the next question is whether they also share a step size. If they did,
the log-bridge would be a full metric identification, not just a structural
one, and CRF would carry a new exact ratio between the coupling-rate scale
and the floor-spacing scale. The numerical probe refuses to cooperate. The
ratio is 36.96, clean to nobody. The finding below records this refusal
carefully, because drawing the boundary of a unification is just as
important as stating it.

%%=============================================================
\section{FIND-ZC88-01: The Steps Are Not Proportional (Gift)}
\label{sec:find-zc88}
%%=============================================================

% [BRIDGE-PROSE: integration-added, Extension Freedom narrative, v3.6 §31.3]
A unification that shows two objects share a shape naturally raises a
sharper question: do they also share a scale? If the two ladder steps were
proportional, the log-bridge would be a metric identification --- the floor
spacing and the rate spacing would be locked together by a single CRF
constant, and the framework would carry a new exact relation between two
quantities derived from different parts of the $\mathbb{Z}_{15}$ machinery.
That would have been a stronger result than what is claimed here.

The numerical probe is unambiguous: the ratio of steps is $36.96\ldots$,
sitting between 36 and 37 with no near-form traceable to CRF structure. In
the tradition of the ZC33 scar and the ZC71 irreducibility result, the
right response is not to manufacture a near-form but to register the
boundary clearly: the unification is structural (shape), not metric (scale).
Naming that boundary is the content of the gift below.

\needspace{8\baselineskip}

\begin{giftbox}[title={FIND-ZC88-01: Structural, Not Metric ---
  $\DeltaGam/(\tfrac14\ln r)$ is Not Clean ($R_0$, Exact --- Gift)}]

\textbf{Statement.}
The bridge of THM-ZC88-01 identifies the \emph{shape} of the two ladders,
not their \emph{step size}. The ratio of steps is
\begin{equation}
  \frac{s_\Gamma}{s_{\rm f}}
  =\frac{\DeltaGam}{\tfrac14\ln r}
  =\frac{\ln 2/\ln 15}{\tfrac14\ln\!\sqrt{(1+2x)/(1+x)}}
  = 36.9638\ldots,
  \label{eq:find-main-zc88}
\end{equation}
which is not a small rational, not $|\mathbb{Z}_{15}|$-clean, and not a
recognised CRF constant. The two ladders share their progression-and-mirror
\emph{form} but run at unrelated rates.

\textbf{Why this is the honest result.}
A ratio of $36.9638\ldots$ sits between $36$ and $37$ with no clean form.
The floor step $\tfrac14\ln r$ lives in the $\mathbb{Q}(\sqrt5)$-flavoured
Born-chain territory of ZC57/67, while $\DeltaGam=\ln2/\ln15$ is a pure
$\mathbb{Z}_{15}$ logarithm; there is no algebraic reason for these to be
commensurate at T-canonical. Forcing a near-integer fit (e.g.\ $36$, or
$|\Ztf|+21$) would repeat the ZC33 error of approximating a structural gap.
The structural identity stands and the metric question is left explicitly open.

\textbf{Classification:} Exact Finding (Gift --- a boundary clearly drawn).
\physmeaning{the two faces of the ladder are genuinely the same shape, but
  the universe did not make their spacings match. The coincidence is
  structural; reading more into the numbers would be the ZC33 mistake again.}
\end{giftbox}

%%=============================================================
\section{Phase Misalignment}
\label{sec:pm-zc88}
%%=============================================================

\begin{derivedbox}[title={PM-ZC88-01: ``The Two Steps Are Proportional''
  --- Probe Refuted (Failure Stance)}]
\textbf{What happened.}
On finding both ladders are APs, the natural probe was: is
$\DeltaGam=\kappa\cdot\tfrac14\ln r$ for a clean CRF constant $\kappa$?

\textbf{Result.}
$\kappa=36.9638\ldots$ --- no clean form. The probe failed to produce a
metric link.

\textbf{Why it was not an error.}
Phase Misalignment, not a wrong turn: the failure \emph{is} the content of
FIND-ZC88-01. It draws the exact line between what the log-bridge does
(unify shape) and does not (unify scale), preventing a future session from
over-claiming a numerical coincidence.

\textbf{Preservation.}
Recorded as FIND-ZC88-01 (boundary) and GAP-ZC88-01 (open question), not
deleted.

\textbf{Pattern type:} tautology test / approximation misclassified
(guards against manufacturing a near-form for $\kappa$).

\textbf{Lesson --- PWA-10:} a clean \emph{shape} match does not license a
clean \emph{metric} match; test the step ratio explicitly before claiming
proportionality.
\end{derivedbox}

%%=============================================================
\section{Open Gap}
\label{sec:gap-zc88}
%%=============================================================

\begin{openqbox}[title={GAP-ZC88-01: Is There a $d_s$-General Exact
  Step-Relation? (Priority~3; ZC89 target)}]
\textbf{Statement.}
At T-canonical ($d_s=3$) the step ratio is $36.9638\ldots$ with no
near-form. Does a closed relation between $s_\Gamma=\DeltaGam$ and
$s_{\rm f}=\tfrac14\ln r$ emerge as an exact function of $d_s$ (and the
$\mathbb{Z}_{15}$ data), of which $36.96$ is merely the $d_s=3$ value?

\textbf{Why this is distinct.}
THM-ZC88-01 and COR-ZC88-01 are exact and closed at $d_s=3$; this gap is
the \emph{metric} question only, deliberately separated from the
\emph{structural} result so the latter is not held hostage to the former.
It also connects to GAP-ZC21-00 (other-$d_s$ behaviour), where the two
$n_m$ conventions diverge (PWA scattered-knowledge note).

\textbf{Approach candidates.}
\begin{enumerate}[leftmargin=2em, noitemsep]
  \item Express both steps as functions of $d_s$ via $x(d_s)$ and the
    sector-weight scheme, then take the ratio symbolically.
  \item Test whether the ratio is $\mathbb{Q}(\sqrt5)$-transcendental
    (cf.\ the ZC71 $0.66\%$ disjointness pattern).
\end{enumerate}
\textbf{Status:} Open. Target: ZC89.
\end{openqbox}

%%=============================================================
\section{Master Status Table}
%%=============================================================

{\small\renewcommand{\arraystretch}{1.3}
\begin{longtable}{@{}p{4.6cm}p{2.4cm}p{5.2cm}@{}}
\toprule
\textbf{Atomic unit} & \textbf{Status} & \textbf{Note} \\
\midrule
\endhead
LEM-ZC88-01 (log-floor AP) & Exact
  & step $\tfrac14\ln r$; axis $\ln\epsuniv$ \\[4pt]
THM-ZC88-01 (co-structure) & Exact
  & both ladders are mirror-symmetric 3-term APs \\[4pt]
COR-ZC88-01 (mirrors = one law) & Exact
  & AM-of-logs $=$ log-of-GM; unifies FIND-ZC81/83-01 \\[4pt]
FIND-ZC88-01 (steps not prop.) & Exact (Gift)
  & $\DeltaGam/(\tfrac14\ln r)=36.96\ldots$; structural not metric \\[4pt]
GAP-ZC88-01 (step relation) & Open
  & Target: ZC89; $d_s$-general step ratio \\[4pt]
PM-ZC88-01 (proportionality) & Scar
  & tautology test / approximation misclassified \\
\midrule
\multicolumn{3}{@{}p{12.5cm}}{\textit{Gap-status changes from prior
  papers:} none. ZC88 closes no prior gap; it unifies two existing Gifts
  (FIND-ZC81-01, FIND-ZC83-01) without altering any value.} \\
\bottomrule
\end{longtable}
}

%%=============================================================
%% ATOMIC REGISTRY (MANDATORY)
%%=============================================================

\begin{formalbox}[title={Atomic Registry --- ZC88}]
\begin{center}
\renewcommand{\arraystretch}{1.3}
\small
\begin{tabular}{lp{8.0cm}l}
\toprule
\textbf{ID} & \textbf{Description} & \textbf{Status} \\
\midrule
LEM-ZC88-01
  & Log of the floor lattice is an AP, step $\tfrac14\ln r$, axis
    $\ln\epsuniv$
  & Exact \\[3pt]
THM-ZC88-01
  & Floor lattice (log) and $\Gamma$-ladder (linear) are co-structural
    mirror-symmetric APs
  & Exact \\[3pt]
COR-ZC88-01
  & FIND-ZC83-01 $\equiv$ FIND-ZC81-01: AM-of-logs $=$ log-of-GM
  & Exact \\[3pt]
FIND-ZC88-01
  & Step ratio $\DeltaGam/(\tfrac14\ln r)=36.96\ldots$ not clean;
    bridge structural not metric
  & Exact (Gift) \\[3pt]
GAP-ZC88-01
  & $d_s$-general exact step-relation between the two ladders; target ZC89
  & Open \\[3pt]
PM-ZC88-01
  & Proportionality probe refuted (pattern: tautology test / approximation
    misclassified)
  & Scar \\
\midrule
\multicolumn{3}{@{}p{11.5cm}}{\textit{Inherited verbatim (not new):}
  LEM-ZC83-01, FIND-ZC83-01 (Part~K); FIND-ZC81-01, THM-ZC73-01,
  FIND-ZC73-01 (Parts~J/I). No value modified.} \\
\bottomrule
\end{tabular}
\end{center}
\end{formalbox}

%%=============================================================
\section{R-Layer Research Note}
\label{sec:rlayer-zc88}
%%=============================================================

\begin{layerbox}[title={R-Layer Assignments for ZC88}]

\textbf{Layer assignments:}

{\small\renewcommand{\arraystretch}{1.25}
\begin{center}
\begin{tabular}{@{}llll@{}}
\toprule
Atomic unit & R-layer & Cost $T$ & Source \\
\midrule
LEM-ZC88-01 & $R_0$ & $0$ & algebraic from LEM-ZC83-01 \\
THM-ZC88-01 & $R_0$ & $0$ & algebraic from THM-ZC73-01 + LEM-ZC88-01 \\
COR-ZC88-01 & $R_0$ & $0$ & log-identity \\
FIND-ZC88-01 & $R_0$ & $0$ & numerical ratio (exact inputs) \\
\bottomrule
\end{tabular}
\end{center}}

\textbf{Two-layer floor (PWA-8):}
\begin{itemize}[leftmargin=2em, noitemsep]
  \item Spectral layer ($R_0$ Wc/Fiedler): $\epscan=0.007550$ ---
    \emph{not used} in this paper.
  \item Bridge layer ($R_0{\to}R_{\max}$): $\epsuniv=0.007599$ ---
    \emph{used} as the floor-lattice axis. Separation $0.66\%$ from
    $\epscan$ irreducible (THM-ZC71-01); ZC88 does not touch it.
\end{itemize}

All ZC88 results are at the $R_0$ layer; no Born-chain bridge running
enters. The $\Gamma$-ladder rungs are likewise $R_0$ values (THM-ZC73-01).
\end{layerbox}

%%=============================================================
\section{Genealogy Map}
\label{sec:genealogy-zc88}
%%=============================================================

\begin{genealogybox}[title={How ZC88's Key Objects Trace Back}]

\textbf{Branch A --- Floor lattice (multiplicative):}
\begin{itemize}[leftmargin=2em, noitemsep]
  \item \textbf{ZC67/THM-ZC67-01}: sector-weight ratio
    $r=\sqrt{(1+2x)/(1+x)}$.
  \item $\to$ \textbf{ZC70}: exact $\epsEM,\epsW$ via $r^{\pm1/4}$.
  \item $\to$ \textbf{ZC83/LEM-ZC83-01, FIND-ZC83-01}: floor formula
    and $\sqrt{\epsEM\epsW}=\epsuniv$ (GM mirror).
  \item $\to$ \textbf{ZC88/LEM-ZC88-01}: logs of these are an AP.
\end{itemize}

\textbf{Branch B --- $\Gamma$-ladder (additive):}
\begin{itemize}[leftmargin=2em, noitemsep]
  \item \textbf{ZC73/THM-ZC73-01}: $\Gamma_k=\Gcoup/\ds-1+k\DeltaGam$,
    step $\DeltaGam=\ln2/\ln15$.
  \item $\to$ \textbf{ZC81/FIND-ZC81-01}: mirror
    $\GammaEM+\GammaS=2\GammaW$ (AM mirror).
  \item $\to$ \textbf{ZC88/THM-ZC88-01}: same AP shape as Branch~A.
\end{itemize}

\textbf{Convergence at ZC88:}
Branch~A is multiplicative and Branch~B is additive; the logarithm carries
A onto B's form. Because the arithmetic mean of logs is the log of the
geometric mean, the GM mirror of Branch~A and the AM mirror of Branch~B
become a single ``middle equals mean'' law (COR-ZC88-01). The convergence
is structural; the step sizes remain unrelated (FIND-ZC88-01).

ZC88's contribution: it shows two separately-discovered mirror symmetries
are one law on two scales, and it draws the exact boundary of that unification.
\end{genealogybox}

%%=============================================================
\section{Closing Sentence}
%%=============================================================
%% TONE B — CONTEMPLATIVE (structural bridge; metric question left open).

\bigskip
\begin{center}
\textit{%
  Two mirrors that never seemed to face each other turned out to be one
  glass, held to the light at two angles.\\[0.3em]
  The shape is settled; the spacing belongs to ZC89.
}

\medskip
{\thaifont
เมื่อมองทั้งระบบพร้อมกัน บันไดที่เคยเป็นสองกลับเป็นหนึ่ง ---
ต่างกันเพียงว่าอ่านด้วยผลคูณหรือผลบวก
ส่วนระยะห่างของขั้น ยังเป็นปริศนาที่รอ ZC89.}
\end{center}

% Governance Note: This document follows Upgrade Protocol v1.3
% and CRF Style Guide v3.6. Gauge: T-canonical. Tone B.


% [BRIDGE-PROSE: integration-added, Style Guide v3.7 §31.6 / UPG-30]
\noindent
ZC88 establishes that the two mirrors are the same law in log-space, but it
deliberately leaves one thing open (GAP-ZC88-01): knowing that a step
\emph{exists} in the additive picture does not say \emph{which physical
layer} that step lives on --- the metric step and the rung step differ by a
factor of two, and the bare theorem cannot say which is the ``real'' one.
Leaving that unfixed would let later papers quietly pick whichever factor
is convenient. ZC89 closes the gap by tying each step to a specific layer.
\section{Background: Completing the ZC88 Arc}
\label{sec:background-zc89}
%%=============================================================

\begin{keybox}[title={If you are just arriving --- ZC89 in plain language}]
\sloppy
ZC88 showed that two patterns in CRF --- a geometric-mean mirror among the
instability \emph{floors} and an arithmetic-mean mirror among the running
\emph{rates} --- are one structure seen on two scales. It then asked a
sharper question and left it open: the two ladders have the same shape, but
do their \emph{step sizes} relate by any clean rule? ZC89 answers yes,
there is an exact formula --- and the formula spits out an ugly irrational
number, $36.96$. That ugliness is the point: the shapes match, the spacings
do not, and now we can prove it in closed form.

The second result is about a floor that does not fit. CRF has four
instability floors. Three of them line up on an evenly spaced ladder. The
fourth, $\epscan$ (the ``canonical'' or spectral floor), sits just off the
ladder --- and the exact size of its offset turns out to equal a gap CRF
already knew was irreducible (the ZC71 two-layer separation). So the ladder
is not a universal law; it is the fingerprint of one specific layer, blind
to the other by design.

\textbf{Minimum vocabulary:}
\textit{floor} ($\varepsilon$) --- minimum sector instability;\;
\textit{rung step} ($\srung$) --- spacing of the bridge ladder in log;\;
\textit{$\Gamma$ step} ($\sGamma$) --- spacing of the rate ladder;\;
\textit{bridge-AP-forbidden} --- not on the ladder.
CRF is a map under construction, not a claim of final truth.
\end{keybox}

\begin{keybox}[title={Numerical Summary --- ZC89 (T-canonical, recomputed)}]
{\small\renewcommand{\arraystretch}{1.25}
\begin{center}
\begin{tabular}{@{}p{6.2cm}p{2.8cm}p{3.4cm}@{}}
\toprule
\textbf{Quantity} & \textbf{Value} & \textbf{Source} \\
\midrule
$x=n\,\epscan$                        & $0.060400$   & inputs \\
$\rborn=(1+2x)/(1+x)$                 & $1.056960$   & ZC67 \\
$\srung=\tfrac12\ln(\epsW/\epsEM)$    & $0.0138834$  & rung (positions) \\
$\smetric=\tfrac14\ln(\epsW/\epsEM)$   & $0.0069246$  & metric (ratio) \\
$\sGamma=\DeltaGam=\ln2/\ln15$        & $0.2559580$  & ZC73 \\
$\sGamma/\smetric$ (closed form)      & $36.9638$    & \textbf{FIND-ZC89-01} \\
$q(\epscan)$ (axis-centred)           & $-0.4660$    & \textbf{THM-ZC89-01} \\
$\deltalayer=|q(\epscan)|$            & $0.4660$     & \textbf{COR-ZC89-01} \\
ZC71 log-gap $/\,\srung$              & $0.4660$     & matches $\deltalayer$ \\
\bottomrule
\end{tabular}
\end{center}}
\end{keybox}

ZC88 ended on a deliberate split. The structural half --- the two ladders
share a mirror-symmetric shape --- was proved and closed. The metric half
--- whether the spacings relate --- was set aside as GAP-ZC88-01, with the
honest note that the numerical ratio $36.96$ showed no clean near-form. The
gap was not a loose end; it was a question framed precisely enough to
answer. This paper answers it, and then turns the same machinery on a
puzzle the corpus had filed separately: the spectral floor $\epscan$ that
never sat comfortably with the other three.

\subsection{What was already known}

\begin{keybox}[title={Pre-ZC89 Status}]
\textbf{Known:}
\begin{itemize}[leftmargin=2em, noitemsep]
  \item ZC88: floors $\epsEM,\epsuniv,\epsW$ form a mirror-symmetric AP in
    log, axis $\ln\epsuniv$, rung step $\srung=\tfrac12\ln(\epsW/\epsEM)$.
  \item ZC88/FIND-ZC88-01: the step ratio $\sGamma/\smetric=36.96\ldots$ is
    not a clean constant (metric, not structural). \textbf{GAP-ZC88-01}
    asked for an exact $d_s$-general expression.
  \item ZC71/THM-ZC71-01: the spectral floor $\epscan$ and the bridge floor
    $\epsuniv$ differ by an irreducible $0.66\%$; they probe different
    $\mathbb{Z}_{15}$ Laplacian invariants. \textbf{KEY PILLAR.}
  \item ZC67: $\rborn=(1+2x)/(1+x)$, $x=n\epscan$ (Born-chain sector ratio).
\end{itemize}
\textbf{Open (before ZC89):} GAP-ZC88-01 (the step relation); and the
unexamined question of \emph{where} $\epscan$ sits relative to the bridge AP.
\end{keybox}

\begin{keybox}[title={What ZC89 Does NOT Do}]
\begin{itemize}[leftmargin=2em, noitemsep]
  \item Does not modify any floor value or any $\Gamma$ rung (all inherited).
  \item Does not claim the step ratio is rational or clean --- it proves the
    opposite, in closed form.
  \item Does not re-open the ZC88 structural results (used as pillars).
  \item Does not assert $\epscan$ belongs on the AP --- it proves exclusion.
\end{itemize}
\end{keybox}

\begin{derivedbox}[title={Discovery Note}]
GAP-ZC88-01 was the metric question, deliberately separated in ZC88.
Reverse tracing (results~$\to$~axioms, per the bidirectional method) showed
the stubborn ratio $36.96$ factors cleanly into $8\ln2/[\ln15\,\ln\rborn]$
once one stops looking for a \emph{rational} and starts looking for a
\emph{structural} expression: each factor is already a named CRF quantity.
The Layer-Specificity result came from the inverse question --- not ``where
is $\epscan$ on the ladder?'' but ``why is it \emph{not}?'' --- and the
answer was a gap the corpus already owned: the ZC71 separation, now read in
units of rung-steps.
\end{derivedbox}

\subsection{Pre-Work Audit}

\begin{formalbox}[title={Pre-Work Audit Record --- ZC89}]
Scanned before writing (PWA-1 through PWA-11):
\begin{itemize}[leftmargin=2em, noitemsep]
  \item ZC88 LEM/THM/COR (bridge AP, rung step) --- \textbf{KEY PILLAR}.
  \item ZC71/THM-ZC71-01 (two-layer irreducibility) --- \textbf{KEY PILLAR}.
  \item ZC67 ($\rborn$), ZC73 ($\DeltaGam=\ln2/\ln15$) --- inherited verbatim.
\end{itemize}
\textbf{PWA-6}: new macros via \texttt{\textbackslash providecommand}; no
Preamble~v3.6 conflict.\\
\textbf{PWA-8} (layer self-consistency): \textbf{this paper is precisely a
PWA-8 statement} --- it identifies which layer each object belongs to and
proves the bridge AP is bridge-layer-specific. $\epscan$ is spectral-layer;
$\epsuniv,\epsEM,\epsW$ are bridge-layer.\\
\textbf{PWA-9} (step-definition): the canonical bridge rung step is fixed as
$\srung=\tfrac12\ln(\epsW/\epsEM)$ (the $\epsEM\!\leftrightarrow\!\epsuniv$
spacing). Prior drafts that used $\tfrac14\ln(\epsW/\epsEM)$ or
$\ln\rborn/8$ are reconciled in \S\ref{sec:stepdef-zc89}.\\
\textbf{PWA-10} (tautology): FIND-ZC89-01 is not circular --- $\srung$ comes
from the floor lattice (Born chain) and $\sGamma$ from the $\Gamma$-ladder
(K14 fixed point); independent sub-chains. See PM-ZC89-PROBE-C-01.
\end{formalbox}

\newpage

%%=============================================================
\section{FIND-ZC89-01: Closed Form for the Step Ratio (closes GAP-ZC88-01)}
\label{sec:find-zc89}
%%=============================================================

Before the formula, the question it answers, stated so an outsider can hold
it: ZC88 found two evenly spaced ladders of the same shape but could not say
whether their step sizes were linked. If they were linked by a clean
constant, the two ladders would be one metric object and CRF would gain a
new exact ratio. If they were not, the match is only a shape coincidence.
The probe in ZC88 returned $36.96$ --- a number clean to nobody --- and the
honest move was to record that and move on. What follows shows that $36.96$,
though irrational, is not arbitrary: it is built entirely from quantities
CRF already derives.

\needspace{8\baselineskip}

\begin{giftbox}[title={FIND-ZC89-01: $\sGamma/\smetric
  = 2^{d_s}\ln2/[\ln(d_s(2^{d_s}\!-\!d_s))\,\ln\rborn]$ ($R_0$, Exact ---
  closes GAP-ZC88-01 and GAP-ZC89-01)}]

\textbf{Statement.}
\begin{equation}
  \boxed{\;
    \frac{\sGamma}{\smetric}
    = \frac{2^{d_s}\ln 2}{\ln\!\big[d_s(2^{d_s}\!-\!d_s)\big]\cdot\ln\rborn}
    \;\xrightarrow{d_s=3}\;
    \frac{8\ln 2}{\ln 15\cdot\ln\rborn}
    = 36.9638\ldots
  \;}
  \label{eq:find-main-zc89}
\end{equation}
with $\rborn=(1+2x)/(1+x)$, $x=n\,\epscan$.

\textbf{Derivation.}

\textit{Step 1 ($\sGamma$).}
From ZC73, $\sGamma=\DeltaGam=\ln 2/\ln 15$.

\textit{Step 2 ($\smetric$, the metric step).}
ZC88 measures the step ratio against $\smetric=\tfrac14\ln(\epsW/\epsEM)$.
The floor lattice (ZC83/ZC67) gives $\epsW/\epsEM=\rborn^{1/2}$ along the
exact $r$-path, so $\ln(\epsW/\epsEM)=\tfrac12\ln\rborn$ and hence
$\smetric=\tfrac18\ln\rborn$.

\textit{Step 3 (ratio at $d_s=3$).}
\[
  \frac{\sGamma}{\smetric}
  = \frac{\ln 2/\ln 15}{\tfrac18\ln\rborn}
  = \frac{8\ln 2}{\ln 15\cdot\ln\rborn}.
\]
Numerically $\ln\rborn=0.055397$, $\ln15=2.70805$, giving $36.9638$.

\textit{Step 4 ($\ln 15$ is a function of $d_s$ --- closes GAP-ZC89-01).}
The group order is not a free input. From the corpus $\delta$-chain:
the Key Identity $3d_s=2^{d_s}+1$ has unique solution $d_s=3$; the matter
split $n=2^{d_s}=d_s+n_m$ gives $n_m=2^{d_s}-d_s$; and CRT identifies the
group order as the product
\[
  |\mathbb{Z}_{15}| = d_s\cdot n_m = d_s\,(2^{d_s}-d_s)
  \;\xrightarrow{d_s=3}\; 3\cdot 5 = 15.
\]
Substituting $\ln 15=\ln[d_s(2^{d_s}-d_s)]$ and $8=2^{d_s}$ yields the
boxed $d_s$-general form. The factor that ZC89 v1 and v2 left as a ``fixed
input'' is therefore fully derived: \textbf{both GAP-ZC88-01 and
GAP-ZC89-01 are closed.}

\textbf{Each factor is $\mathbb{Z}_{15}$-structural and $d_s$-derived:}
\begin{center}
{\small\renewcommand{\arraystretch}{1.2}
\begin{tabular}{@{}lll@{}}
\toprule
Factor & Meaning & $d_s$-form \\
\midrule
$8$         & ladder normalisation        & $2^{d_s}$ \\
$\ln 2$     & $\Gamma$-ladder bit (ZC73)   & $\ln 2$ \\
$\ln 15$    & group order (CRT)            & $\ln[d_s(2^{d_s}\!-\!d_s)]$ \\
$\ln\rborn$ & Born-chain weight (ZC67)     & $\ln\rborn$ \\
\bottomrule
\end{tabular}}
\end{center}

\textbf{Status:} Exact Finding (Gift). Closed-form and $d_s$-general;
irrational, no near-form --- the metric content of FIND-ZC88-01, now fully
derived. \hfill$\square$
\physmeaning{the two ladders are built from the same alphabet ($2$, $d_s$,
  $n_m$, $\rborn$) but combined differently --- additively for $\Gamma$,
  geometrically for the floors. Their step ratio is a ratio of a logarithm
  to a logarithm, generically irrational. The ugliness of $36.96$ is
  structural honesty, not noise --- and every symbol in it traces to $d_s$
  through one Key Identity.}
\end{giftbox}

\begin{derivedbox}[title={On the rounding artefact (honest note)}]
Computed from the published 6-significant-figure floors, the naive ratio
differs from $36.9638$ by about $0.25\%$, and the rung positions of
$\epsEM,\epsW$ come out as $-1.0022,+0.9978$ rather than exactly
$\mp 1$. This is a rounding artefact: the ZC88 AP is exact only along the
algebraic $r$-path, where $\epsW/\epsEM=\rborn^{1/2}$ holds identically. The
closed form \eqref{eq:find-main-zc89} uses that exact path and is therefore the
correct value; the published decimals are a lossy display of it. No result
here rests on the third decimal of a floor.
\end{derivedbox}

%%=============================================================
\section{The Canonical Bridge Step (reconciliation)}
\label{sec:stepdef-zc89}
%%=============================================================

Two log-step units appear in this arc, and conflating them is exactly the
trap that produced an error in an earlier draft (recorded as
PM-ZC89-FACTOR2-01 below). They differ by a factor of two and serve
different purposes; both are pinned here.

\begin{formalbox}[title={Two bridge steps, pinned (PWA-9 lock)}]
Let $\rfloor:=\epsW/\epsEM=\rborn^{1/2}$ (exact $r$-path, LEM-ZC83-01).
\begin{itemize}[leftmargin=2em, noitemsep]
  \item \textbf{Metric step} (used in the ZC88/ZC89 step \emph{ratio}):
    \[
      \smetric := \tfrac14\ln\rfloor = \tfrac18\ln\rborn = 0.00692457.
    \]
    This is ZC88's $s$; $\sGamma/\smetric=36.9638$.
  \item \textbf{Rung step} (the spacing between adjacent AP floors, used for
    \emph{positions}):
    \[
      \srung := \ln\epsuniv-\ln\epsEM = \tfrac12\ln\rfloor
              = \tfrac14\ln\rborn = 2\,\smetric = 0.0138834.
    \]
\end{itemize}
The relation is simply $\srung=2\smetric$: the EM--univ \emph{rung} is two
metric steps, because ZC88 placed the axis at $\ln\epsuniv$ and measured
displacement to a floor as $\smetric$, half the rung. The draft form
$\ln\rborn/8$ equals $\smetric$; the form $\tfrac12\ln(\epsW/\epsEM)$
equals $\srung$. No floor value changes; only the two units are now
named distinctly so the ratio (in $\smetric$) and the positions (in
$\srung$) never cross wires again.
\end{formalbox}

%%=============================================================
\section{THM-ZC89-01: The Layer-Specificity Theorem}
\label{sec:thm-zc89}
%%=============================================================

The first half of this paper closed a metric question. The second half uses
the same coordinate to settle a structural one that the corpus had left
implicit. CRF has four floors, not three: alongside the bridge trio
$\epsEM,\epsuniv,\epsW$ there is the spectral floor $\epscan$, defined by a
different fixed-point equation (ZC79). It has always sat slightly apart ---
$0.66\%$ from $\epsuniv$, an offset ZC71 proved irreducible. The natural
question is whether $\epscan$ nonetheless lands on the bridge ladder. It
does not, and the reason is exact.

\needspace{10\baselineskip}

\begin{formalbox}[title={THM-ZC89-01: Layer-Specificity ---
  $\epscan$ is Bridge-AP-Forbidden ($R_0$, Exact)}]

\textbf{Statement.}
Define the axis-centred bridge coordinate of any floor value, with axis
$\ln\epsuniv$ and unit one rung step $\srung$:
\begin{equation}
  q(\varepsilon) := \frac{\ln\varepsilon-\ln\epsuniv}{\srung}.
  \label{eq:position-zc89}
\end{equation}
Then the bridge floors sit on integer rungs, and $\epscan$ does not:
\begin{equation}
  \boxed{\;
    q(\epsEM)=-1,\quad q(\epsuniv)=0,\quad q(\epsW)=+1,
    \qquad
    q(\epscan)=-\deltalayer\notin\mathbb{Z},
  \;}
  \label{eq:thm-main-zc89}
\end{equation}
with
\begin{equation}
  \deltalayer
  := \frac{\ln\epsuniv-\ln\epscan}{\srung}
   = 0.4660 \;>\;0.
  \label{eq:delta-zc89}
\end{equation}

\textbf{Proof.}

\textit{Step 1 (bridge floors are integer rungs).}
By definition $q(\epsuniv)=0$. Since $\srung=\ln\epsuniv-\ln\epsEM$,
$q(\epsEM)=(\ln\epsEM-\ln\epsuniv)/\srung=-1$, and by mirror symmetry
(ZC88) $q(\epsW)=+1$. (With published decimals these read $-1.0022$,
$+0.9978$; exact along the $r$-path --- see \S\ref{sec:find-zc89} note.)

\textit{Step 2 (position of $\epscan$).}
Directly from \eqref{eq:position-zc89} and \eqref{eq:delta-zc89},
\[
  q(\epscan)
  = \frac{\ln\epscan-\ln\epsuniv}{\srung}
  = -\frac{\ln\epsuniv-\ln\epscan}{\srung}
  = -\deltalayer
  = -0.4660.
\]
Because $\epscan<\epsuniv$, this is \emph{negative}: $\epscan$ lies below
the axis, between rung $0$ ($\epsuniv$) and rung $-1$ ($\epsEM$). Its
distance to the nearer rung is $\min(0.466,0.534)=0.466$, far larger than
the rounding noise ($\sim 0.002$).

\textit{Step 3 ($\deltalayer\neq$ integer, exclusion).}
$\deltalayer=0$ would require $\epscan=\epsuniv$; $\deltalayer=1$ would
require $\epscan=\epsEM$. THM-ZC71-01 forbids both: the spectral
fixed-point (probing $\lambda_2(\mathbb{Z}_{15})^{-1}$) and the bridge
fixed-point (probing $\lambda_2^2=n_m/\varphi^2$) are different invariants,
so $\ln\epsuniv-\ln\epscan\neq 0$ irreducibly, and its size ($0.66\%$ in
$\varepsilon$, $0.466$ in rungs) is fixed, not free. Hence $q(\epscan)$ is
non-integer and $\epscan$ cannot lie on the bridge AP. \hfill$\square$

\textbf{Status:} Exact Theorem.
\physmeaning{the bridge ladder is not a law every floor must obey. It is the
  fingerprint of the bridge layer's Born-chain geometry. $\epscan$ belongs
  to a different layer (spectral), and asking it to stand on the bridge
  ladder is asking one $\mathbb{Z}_{15}$ invariant to equal another --- the
  exact thing ZC71 proved impossible.}
\end{formalbox}

%%=============================================================
\section{COR-ZC89-01: The Displacement Equals the ZC71 Gap}
\label{sec:cor-zc89}
%%=============================================================

\begin{formalbox}[title={COR-ZC89-01: $\deltalayer$ is the ZC71 gap in
  rung units ($R_0$, Exact)}]
\textbf{Statement.}
\begin{equation}
  \deltalayer
  = \frac{\ln\epsuniv-\ln\epscan}{\srung}
  = \frac{(\text{ZC71 two-layer log-gap})}{\srung}
  = 0.4660.
\end{equation}
\textbf{Resolution.} The numerator is, verbatim, the ZC71 two-layer
separation; dividing by the rung step expresses it in the bridge ladder's
own units. So the amount by which $\epscan$ misses the ladder is not a new
free quantity --- it is the irreducible ZC71 gap, re-measured. The two
results of this paper meet here: the ladder's spacing ($\srung$, via
FIND-ZC89-01) sets the ruler, and the ZC71 gap sets the distance.
\textbf{Status:} Exact Corollary.
\end{formalbox}

%%=============================================================
\section{Phase Misalignment}
\label{sec:pm-zc89}
%%=============================================================

\begin{derivedbox}[title={PM-ZC89-FACTOR2-01: The Factor-of-Two
  Unit Conflation --- Caught in Review}]
\textbf{What happened.}
A draft of this paper (v2) defined the step used in the ratio as
$\srung=\tfrac12\ln(\epsW/\epsEM)$ in one section while the abstract and
giftbox used the ZC88 form $8\ln2/[\ln15\ln\rborn]$. The two are
inconsistent by a factor of two: the boxed formula evaluates to $36.96$,
but the section-locked $\srung$ gives $\sGamma/\srung=18.48$. The paper
silently contradicted itself.

\textbf{Root cause.}
Two distinct quantities both felt like ``the step'': the metric step
$\smetric=\tfrac14\ln(\epsW/\epsEM)$ that ZC88 measures the ratio against,
and the rung step $\srung=\tfrac12\ln(\epsW/\epsEM)=2\smetric$ that is the
actual spacing between adjacent floors. Compounding it, two different
$r$'s carry similar names: $\rfloor=\epsW/\epsEM=1.0282$ and
$\rborn=1.0570$, with $\rfloor=\rborn^{1/2}$. Mixing $\ln\rfloor$ with
$\ln\rborn$ moved factors of two around invisibly.

\textbf{Why it is preserved, not erased.}
The error is instructive: it shows the precise failure mode of an
\emph{undeclared unit}. A formula can be individually correct in two places
and jointly contradictory if the two places measure in different units
without saying so. The fix was not new mathematics --- it was naming
$\smetric$ and $\srung$ as separate symbols (\S\ref{sec:stepdef-zc89}) so the
ratio and the positions can never again be computed against silently
different rulers.

\textbf{Pattern type:} undeclared-unit conflation; same word (``step'')
bound to two values differing by a constant factor.

\textbf{Lesson --- PWA-9:} when two quantities share a name and differ by a
clean factor, that cleanliness hides the bug. Give each its own symbol
\emph{before} writing any ratio. A factor of two is the easiest error to
make and the hardest to see, precisely because $2$ looks like it could be a
real structural constant.
\end{derivedbox}

\begin{derivedbox}[title={PM-ZC89-PROBE-C-01: Cross-Domain Probe Refuted}]
\textbf{What happened.}
A probe asked whether mixing floors and rates across domains,
$\epsEM e^{\GammaEM}+\epsW e^{\GammaW}$, reproduces a $\Gamma$ value ---
a tempting ``unification'' of the two ladders into one expression.

\textbf{Result (recomputed).}
$\epsEM e^{\GammaEM}+\epsW e^{\GammaW}=0.013644$, while $\GammaEM=0.013664$.
The two agree to $0.1\%$.

\textbf{Why it is still refuted.}
Numerical proximity is not derivation. There is no structural reason to add
a floor times $e^{\Gamma}$ across two different sectors; the expression
mixes a Born-chain quantity with a running rate with no shared invariant.
A $0.1\%$ coincidence among small numbers is weak evidence, and ZC89 has
just shown (THM-ZC89-01) that the two ladders are layer-distinct ---
precisely the structure that would forbid such a naive cross-term from
being fundamental. The probe is recorded, not promoted.

\textbf{Correction note.}
An earlier draft reported this probe as ``$=0.5027$, gap $86\%$.'' Those
numbers were miscomputed; the true sum is $0.013644$ and the true gap is
$0.1\%$. The refutation does \emph{not} depend on a large gap --- it depends
on the absence of structural motivation. The corrected small gap actually
sharpens the lesson: a near-miss is the most seductive kind of false
pattern, and the guard against it is derivation, not distance.

\textbf{Pattern type:} cross-domain term with no shared invariant;
near-miss mistaken for identity.

\textbf{Lesson --- PWA-10:} proximity is not proof. Demand a derivation path
before treating a numerical coincidence as structure --- especially when it
is close.
\end{derivedbox}

%%=============================================================
\section{Master Status Table}
%%=============================================================

{\small\renewcommand{\arraystretch}{1.3}
\begin{longtable}{@{}p{4.8cm}p{2.2cm}p{5.2cm}@{}}
\toprule
\textbf{Atomic unit} & \textbf{Status} & \textbf{Note} \\
\midrule
\endhead
FIND-ZC89-01 (step ratio) & Exact (Gift)
  & $8\ln2/[\ln15\ln\rborn]=36.9638$; \textbf{closes GAP-ZC88-01} \\[4pt]
THM-ZC89-01 (layer-specificity) & Exact
  & $\epscan$ at $q=-0.4660$, bridge-AP-forbidden \\[4pt]
COR-ZC89-01 (displacement) & Exact
  & $\deltalayer=$ ZC71 gap $/\srung=0.4660$ \\[4pt]
$\srung$ canonical lock (\S\ref{sec:stepdef-zc89}) & Exact
  & $\srung=\tfrac12\ln(\epsW/\epsEM)=\tfrac14\ln\rborn$ \\[4pt]
PM-ZC89-PROBE-C-01 & Scar
  & cross-domain near-miss; refuted (no motivation) \\[4pt]
\midrule
\multicolumn{3}{@{}p{12.5cm}}{\textit{Gap-status changes:}
  \textbf{GAP-ZC88-01 and GAP-ZC89-01 both CLOSED} by FIND-ZC89-01: the step ratio is $d_s$-general closed form, and $\ln15=\ln[d_s(2^{d_s}\!-\!d_s)]$ is pinned by the Key Identity.} \\
\bottomrule
\end{longtable}}

\begin{formalbox}[title={Atomic Registry --- ZC89}]
\begin{center}
\renewcommand{\arraystretch}{1.3}\small
\begin{tabular}{lp{7.6cm}l}
\toprule
\textbf{ID} & \textbf{Description} & \textbf{Status} \\
\midrule
FIND-ZC89-01 & Closed form $\sGamma/\smetric=8\ln2/[\ln15\ln\rborn]$;
  closes GAP-ZC88-01 \& GAP-ZC89-01 & Exact (Gift) \\[3pt]
THM-ZC89-01 & $\epscan$ bridge-AP-forbidden; $q(\epscan)=-0.4660$
  non-integer & Exact \\[3pt]
COR-ZC89-01 & $\deltalayer=$ ZC71 log-gap $/\srung=0.4660$ & Exact \\[3pt]
PM-ZC89-PROBE-C-01 & cross-domain probe refuted (near-miss, no motivation)
  & Scar \\[3pt]
GAP-ZC89-01 & $\ln15=\ln[d_s(2^{d_s}\!-\!d_s)]$ via Key Identity; group order is $d_s$-derived
  & \textbf{Closed} \\
\midrule
\multicolumn{3}{@{}p{11.5cm}}{\textit{Inherited verbatim:} ZC88 bridge AP
  + rung step; ZC71 two-layer gap; ZC67 $\rborn$; ZC73 $\DeltaGam$.
  No value modified.} \\
\bottomrule
\end{tabular}
\end{center}
\end{formalbox}

%%=============================================================
\section{R-Layer Research Note}
\label{sec:rlayer-zc89}
%%=============================================================

\begin{layerbox}[title={R-Layer Assignments for ZC89}]
{\small\renewcommand{\arraystretch}{1.25}
\begin{center}
\begin{tabular}{@{}llll@{}}
\toprule
Atomic unit & R-layer & Cost $T$ & Source \\
\midrule
FIND-ZC89-01 & $R_0$ & $0$ & algebraic (ZC73 + ZC88 rung) \\
THM-ZC89-01 & $R_0$ & $0$ & ZC88 AP + ZC71 gap \\
COR-ZC89-01 & $R_0$ & $0$ & restatement in rung units \\
\bottomrule
\end{tabular}
\end{center}}

\textbf{This paper is itself a layer statement.} Its content is that the
bridge AP is bridge-layer-specific: $\epsEM,\epsuniv,\epsW$ are bridge-layer
($R_0\!\to\!R_{\max}$ Born chain); $\epscan$ is spectral-layer ($R_0$
Wc/Fiedler). THM-ZC89-01 proves the AP cannot cross between them, which is
the geometric face of THM-ZC71-01's irreducibility.
\end{layerbox}

%%=============================================================
\section{Genealogy Map}
\label{sec:genealogy-zc89}
%%=============================================================

\begin{genealogybox}[title={How ZC89's Results Trace Back}]
\textbf{FIND-ZC89-01 (step ratio):}
\begin{itemize}[leftmargin=2em, noitemsep]
  \item ZC73 $\to$ $\sGamma=\ln2/\ln15$.
  \item ZC67 $\to$ $\rborn$; ZC83/ZC88 $\to$ $\smetric=\tfrac18\ln\rborn$ (metric step).
  \item $\to$ ratio $=2^{d_s}\ln2/[\ln(d_s(2^{d_s}\!-\!d_s))\ln\rborn]$, closing GAP-ZC88-01 \& GAP-ZC89-01.
\end{itemize}
\textbf{THM-ZC89-01 (layer-specificity):}
\begin{itemize}[leftmargin=2em, noitemsep]
  \item ZC88 $\to$ bridge AP, axis $\ln\epsuniv$, rung $\srung$.
  \item ZC71 $\to$ irreducible $\epscan$--$\epsuniv$ gap.
  \item $\to$ $q(\epscan)=-(\text{gap})/\srung=-0.4660$, non-integer
    $\Rightarrow$ forbidden.
\end{itemize}
\textbf{Convergence.} FIND sets the ladder's spacing; THM uses that spacing
as a ruler to measure how far the spectral floor misses it; COR identifies
that distance as the ZC71 gap itself. One coordinate, two questions, both
closed.
\end{genealogybox}

%%=============================================================
\section{Closing Sentence}
%%=============================================================
%% TONE B — CONTEMPLATIVE (a gap closed; a layer boundary made visible).

\bigskip
\begin{center}
\textit{%
  The spacing that ZC88 left open turned out to be spelled in the same four
  letters as everything else --- $2$, $15$, $\rborn$, and the dimension
  three.\\[0.3em]
  And the floor that would not join the ladder was not broken; it was
  simply standing in another room.
}

\medskip
{\thaifont
ระยะห่างที่ ZC88 เว้นไว้ กลายเป็นคำที่สะกดด้วยอักษรเดิมสี่ตัว ---
สอง สิบห้า ค่าน้ำหนัก และมิติสาม\\
ส่วนพื้นที่ไม่ยอมขึ้นบันได มิได้ชำรุด --- เพียงแต่ยืนอยู่คนละห้องเท่านั้น.}
\end{center}

% Governance Note: Upgrade Protocol v1.3, CRF Style Guide v3.7. Tone B.


%%=============================================================
\part{Mind-Layer Structure: Quantum-Pair, Irrationality, Barami, Gate (v17.16)}
\label{part:LVIII}
%%=============================================================

% [BRIDGE-PROSE: integration-added, Style Guide v3.7 §31.6 / UPG-30]
\noindent
Cluster 1 worked entirely in the force layer. The mind layer, built in
Part F from the four-way decision structure ($\mathcal{C}_4$) and the
Barami pillars, has its own pair of steps: a smooth, continuous adjustment
(\textit{visesa}) and a sudden, discontinuous one (\textit{nibbedha}). The
open question is whether the discrete/continuous distinction that organised
the force layer is a coincidence of that layer or a structural feature that
repeats. ZC90 imports the pair explicitly and asks whether the two mind-steps
form the same kind of quantum-pair object the force layer showed.
\section{Background: Two Layers That Were Saying One Thing}
\label{sec:background-zc90}
%%=============================================================

%% ── [v3.6 §31.2 / UPG-25] READER'S ON-RAMP ───────
CRF is a single-primitive framework: from one undefined act of
distinction ($\delta$) it reconstructs structures that, in many places,
match the constants of physics. Two of those structures live on very
different ``layers.'' One is the \emph{force layer}: the ladder of
interaction strengths (electromagnetic, weak, and so on) that recent
papers showed to be rungs of a single arithmetic ladder. The other is
the \emph{mind layer}: the description of how a conscious agent moves,
moment by moment, either toward or away from a stable ground. This paper
makes a single observation that links them. On both layers, motion comes
in exactly two flavours --- a small smooth step and a sudden decisive
leap --- and on both layers those two flavours are provably impossible to
express as multiples of each other. The technical sections below make
this precise; the claim itself needs no equation to state: \emph{nature
writes the same two-flavour rule once for forces and once for minds.}

% [BRIDGE-PROSE: integration-added not required here — this is standalone-origin
%  on-ramp prose per UPG-25; no marker needed for in-paper Background.]

%% ── [NEW BLOCK 1] NUMERICAL SUMMARY ────────────────────────
\begin{keybox}[title={Numerical Summary --- ZC90 (T-canonical)}]
{\footnotesize\renewcommand{\arraystretch}{1.3}
\begin{center}
\begin{tabular}{@{}p{5.2cm}p{3.0cm}p{4.0cm}@{}}
\toprule
\textbf{Quantity} & \textbf{Value} & \textbf{Source / Status} \\
\midrule
$\sGam/\srung$ (force ratio)
  & $36.9638\ldots$
  & FIND-ZC89-01 (irrational) \\
$\srung=\tfrac12\ln(\epsW/\epsEM)$
  & $0.01388338$
  & ZC89 (bridge AP step) \\
$\Delta\Gamma=\ln2/\ln15$
  & $0.25595802$
  & FIND-ZC73-01 \\
$\dv$ continuous step (\emph{visesa})
  & $>0$, mag.\ open
  & EQ-MB01 (Part~F) \\
$\db$ discont.\ step (\emph{nibbedha})
  & $\db\gg\dv$
  & EQ-MB01 / DEF-MB03 \\
$\db/\dv$ (mind ratio)
  & \textbf{undetermined}
  & \textbf{GAP-ZC90-01 (new)} \\
\bottomrule
\end{tabular}
\end{center}}
\end{keybox}

%% ── PHYSICS CONTEXT ─────────────────────────────────────────
\subsection{What Each Layer Already Established}

The force-layer fact is recent and sharp. ZC88 (the Log-Bridge Theorem)
showed that the multiplicative floor lattice and the additive Gamma
ladder are one arithmetic progression seen on two scales, joined by the
logarithm. ZC88 left one honest gap: the two ladders are co-structural,
but their \emph{steps} are not obviously proportional, and the measured
ratio $36.9638\ldots$ did not look ``clean.'' ZC89 closed that gap by
proving the ratio in closed form,
\[
  \frac{\sGam}{\srung}=\frac{8\ln 2}{\ln 15\cdot\ln\rborn}=36.9638\ldots,
\]
with every factor a $\Ztf$-structural quantity, and --- crucially ---
proving the ratio \emph{irrational}. The force layer therefore carries
two genuine, co-existing steps that cannot be collapsed into one.

The mind-layer fact is old and was stated in words, not symbols. DEF-MB03
(Part~F, Mind/Barami Companion) classifies the agentive choice into four
movements. Two of them are \emph{forward} --- toward the independent
ground $\Sind$: \emph{visesa-bhāgiya} (\textbf{วิเสสภาคิยะ}), an ordinary
advance, and \emph{nibbedha-bhāgiya} (\textbf{นิพเพธภาคิยะ}), a
penetrative breakthrough. DEF-MB03 states explicitly that the distinction
between them ``is not a difference of magnitude but of kind,'' and that
the four-way choice ``cannot be reduced to a continuous variable.''
EQ-MB01 assigns them KL-divergence increments $-\dv$ and $-\db$ with
$\db\gg\dv$. Read together, these say: one smooth forward step, one
discontinuous forward step, of different kind. That is the same shape as
the force layer --- which is the entire content of this paper.

\subsection{What Was Already Known Before ZC90}

\begin{keybox}[title={Pre-ZC90 Status}]
\textbf{Known:}
\begin{itemize}[leftmargin=2em, noitemsep]
  \item Force layer: $\sGam/\srung=36.9638\ldots$ is exact and
    \emph{irrational} (FIND-ZC89-01, Part~K-adjacent standalone).
  \item Force layer: continuous step $\srung$ and discontinuous
    floor/Gamma step co-exist as a bridge AP (THM-ZC88-01).
  \item Mind layer: \emph{visesa}/\emph{nibbedha} differ ``in kind not
    magnitude''; $\mathcal{C}_4$ is irreducible to a continuous variable
    (DEF-MB03).
  \item Mind layer: Barami is defined only on
    $\Cfourplus=\{\text{visesa},\text{nibbedha}\}$ (DEF-MB04).
\end{itemize}

\textbf{Open:}
\begin{itemize}[leftmargin=2em, noitemsep]
  \item Whether the two layers' two-step structures are the \emph{same}
    structure or merely a loose analogy.
  \item What algebraic object the four-choice $\mathcal{C}_4$ actually is
    (a probe had proposed Klein-four $\mathbb{Z}_2\times\mathbb{Z}_2$).
\end{itemize}

ZC90 answers the first by defining the shared signature (Q-pair) and
proving both layers carry it; it answers the second negatively (the
Klein-four reading fails, PM-ZC90-01) and extracts the correct object
(FIND-ZC90-01). The numeric mind-layer ratio is deferred to GAP-ZC90-01.
\end{keybox}

%% ── [NEW BLOCK 2] SCOPE EXCLUSION ───────────────────────────
\begin{keybox}[title={What ZC90 Does NOT Do}]
\begin{itemize}[leftmargin=2em, noitemsep]
  \item Does not assign a numeric value to $\db/\dv$ (GAP-ZC90-01);
    no calibration of mind-layer step magnitudes exists yet.
  \item Does not modify any $\varepsilon$ floor, $\Gamma$ rung, or any
    ZC88/ZC89 result (used as pillars, unchanged).
  \item Does not modify DEF-MB01--07 or EQ-MB01--04 (used verbatim).
  \item Does not claim Barami is the \emph{same} charge as a force-sector
    charge --- only that both carry the same \emph{structural} signature.
  \item Does not revive the Klein-four reading (refuted; PM-ZC90-01).
\end{itemize}
\end{keybox}

%% ── [NEW BLOCK 3] COMPANION PAPERS ─────────────────────────
\begin{keybox}[title={Companion Papers --- ZC90}]
{\footnotesize\renewcommand{\arraystretch}{1.3}
\begin{center}
\begin{tabular}{@{}p{5.6cm}p{6.9cm}@{}}
\toprule
\textbf{Paper} & \textbf{What it provides to ZC90} \\
\midrule
\texttt{CRF\_ZC89\_LayerSpecificity\_v3.tex}
  & FIND-ZC89-01: force-layer ratio $36.9638$, irrational; $\srung$
    definition --- \textbf{KEY PILLAR (force layer)} \\
\texttt{CRF\_ZC88\_LogBridge\_v2.tex}
  & THM/LEM/COR-ZC88-01: the co-AP bridge with step $\srung$ \\
\texttt{CRF\_Complete\_v17\_16\_Part\_F.tex}
  & DEF-MB03/04, EQ-MB01: agentive $\mathcal{C}_4$, Barami ---
    \textbf{KEY PILLAR (mind layer)} \\
\texttt{CRF\_Complete\_Index\_v17\_16.tex}
  & Part~F (Mind/Barami), Part~K (ZC83--87) context \\
\bottomrule
\end{tabular}
\end{center}}
\end{keybox}

%% ── DISCOVERY NOTE (promoted: the origin was a refuted probe) ──
\begin{derivedbox}[title={Discovery Note}]
This paper began with a probe that asked whether the four-way agentive
choice $\mathcal{C}_4$ was a Klein-four group $\mathbb{Z}_2\times\mathbb{Z}_2$. The
probe was refuted against the corpus: the ``agency'' axis is 3-vs-1, not
2-vs-2 (the \emph{ṭhiti} branch is a passive \emph{anicca} drift, not a
directed choice), and DEF-MB04 singles out the forward pair, which no
symmetric Klein subgroup does. The proof direction became clear only when
the refutation was read the other way: the asymmetry that broke Klein was
itself the signal. The two forward choices are not cosets --- they are a
continuous and a discontinuous quantum of one conserved current, the same
object ZC89 had just found at the force layer. The bottom-up refutation
preceded the top-down theorem.
\end{derivedbox}

%% ── PRE-WORK AUDIT ──────────────────────────────────────────
\subsection{Pre-Work Audit (PWA-1 through PWA-10)}

\begin{formalbox}[title={Pre-Work Audit Record --- ZC90}]
Scanned before writing (Upgrade Protocol v1.3, PWA-1 through PWA-10):
\begin{itemize}[leftmargin=2em, noitemsep]
  \item ZC89/FIND-ZC89-01: force-layer irrational ratio --- \textbf{KEY PILLAR}.
  \item ZC88/THM-ZC88-01: co-AP bridge, step $\srung$ --- \textbf{KEY PILLAR}.
  \item Part~F/DEF-MB03, DEF-MB04, EQ-MB01: agentive $\mathcal{C}_4$,
    Barami --- \textbf{KEY PILLAR}. Read verbatim, not reconstructed.
  \item Part~F/DEF-MB04: Barami defined on $\Cfourplus$ --- the fact that
    refutes Klein-four and motivates FIND-ZC90-01.
\end{itemize}

\textbf{PWA-1} (adjacent-paper read): ZC88, ZC89, and the Part~F MB
cluster read in full before writing; no machinery duplicated.\\
\textbf{PWA-6} (macro convention): paper-specific macros
($\srung,\sGam,\rborn,\dv,\db,\Cfourplus,\Qpair$) all
\texttt{\textbackslash providecommand}; none collide with Preamble v3.1.\\
\textbf{PWA-7} (sector-floor): not triggered (no R-layer bridge computed;
$\epsEM,\epsW$ enter only inside the inherited ZC89 ratio).\\
\textbf{PWA-8} (layer self-consistency): not triggered (no new
$\varepsilon$ value introduced).\\
\textbf{PWA-9} (spectral convention): not triggered.\\
\textbf{PWA-10} (tautology audit): \emph{actively triggered and recorded}
--- the Klein-four probe tested a definition, not physics; logged as
PM-ZC90-01 (pattern: tautology test).\\
\textbf{No duplicate atomic units found.}
\end{formalbox}

\newpage

%%=============================================================
%% ATOMIC UNITS — DEF → LEM → THM → COR → FIND → CONJ → GAP → PM
%%=============================================================

%%-------------------------------------------------------------
\section{The Quantum-Pair Structure\quad(DEF-ZC90-01)}
\label{sec:def-zc90}
%%-------------------------------------------------------------

% [BRIDGE-PROSE: stakes before the box, UPG-30 prose floor]
Before stating either layer's case we need a single name for the thing
both layers turn out to have. The temptation is to reach for a group ---
a $\mathbb{Z}_2$ grading, a cyclic ladder --- but a group is the wrong tool,
because the defining feature here is precisely that the two pieces
\emph{cannot} be composed into a third of the same kind. What we need is
weaker than a group and sharper than an analogy: an ordered direction
along which two generators act, one continuously and one discontinuously,
with the explicit demand that no integer or rational relation ties them
together. We call it a Quantum-Pair.

\needspace{8\baselineskip}
\begin{formalbox}[title={DEF-ZC90-01: Quantum-Pair (Q-pair) ($R_0$, Definition)}]

\textbf{Statement.}
A \emph{Quantum-Pair} $\Qpair=(\mathcal{D};\,g_c,g_d)$ consists of
\begin{enumerate}[leftmargin=2em, noitemsep]
  \item a \emph{conserved direction} $\mathcal{D}$: a signed axis along
    which a quantity accumulates monotonically under the pair's action;
  \item a \emph{continuous generator} $g_c$ producing increments of
    arbitrarily small magnitude $\mu_c>0$ along $\mathcal{D}$;
  \item a \emph{discontinuous generator} $g_d$ producing a fixed increment
    of magnitude $\mu_d$ along the same $\mathcal{D}$, with $\mu_d\gg\mu_c$;
\end{enumerate}
subject to the \textbf{non-commensurability condition}
\begin{equation}
  \boxed{\;\frac{\mu_d}{\mu_c}\notin\mathbb{Q}\quad\text{(or, weakly:
  no rational relation realises $g_d$ as a finite iterate of $g_c$).}\;}
  \label{eq:qpair-zc90}
\end{equation}

\textbf{Remark (why not a group).} A Q-pair is deliberately \emph{not}
closed: $g_c$ and $g_d$ do not generate a third generator of comparable
status, and there is no inverse demanded. The pair is \emph{pointed} (the
direction $\mathcal{D}$ is oriented) and \emph{graded by kind}
(continuous vs discontinuous), not by magnitude. This is exactly the
content DEF-MB03 expresses as ``of kind, not of magnitude.''

\textbf{Status:} Definition.
\end{formalbox}

%%-------------------------------------------------------------
\section{The Force Layer Is a Q-Pair\quad(LEM-ZC90-01)}
\label{sec:lem01-zc90}
%%-------------------------------------------------------------

% [BRIDGE-PROSE: stakes before the box]
The force layer supplies the cleaner of the two instantiations, because
ZC89 already did the hard work: it not only exhibited the two steps but
proved their ratio irrational. All that remains is to read ZC89's result
in the language of Definition~\ref{sec:def-zc90} and confirm each clause holds.

\needspace{8\baselineskip}
\begin{formalbox}[title={LEM-ZC90-01: The Force-Layer Pair Is a Q-Pair
  ($R_0$, Exact)}]

\textbf{Statement.} The pair $(\,\mathcal{D}_{\rm force};\,\srung,\sGam)$
is a Q-pair, where $\mathcal{D}_{\rm force}$ is the log-bridge axis of
THM-ZC88-01, $\srung$ the continuous bridge-AP step, and $\sGam$ the
Gamma-ladder step.

\textbf{Derivation.}

\textit{Step 1 (conserved direction).}
THM-ZC88-01 places the floor lattice and the Gamma ladder on one
arithmetic progression under the log-map; the shared axis (oriented
toward increasing $\Gamma$) is the conserved direction $\mathcal{D}_{\rm
force}$. Both steps accumulate monotonically along it.

\textit{Step 2 (two generators of different kind).}
$\srung=\tfrac12\ln(\epsW/\epsEM)=0.01388338$ is the small continuous
step of the bridge AP; $\sGam=\Delta\Gamma$-scale Gamma step is the large
step. Numerically $\sGam/\srung=36.9638$, so $\mu_d\gg\mu_c$ holds.

\textit{Step 3 (non-commensurability).}
FIND-ZC89-01 proved
$\sGam/\srung=8\ln2/[\ln15\cdot\ln\rborn]$ is \emph{irrational}. Hence
$\mu_d/\mu_c\notin\mathbb{Q}$ exactly, and condition~\eqref{eq:qpair-zc90}
holds in its strong form.

\hfill$\square$

\textbf{Numerical verification (T-canonical):}
\begin{center}
{\small\renewcommand{\arraystretch}{1.25}
\begin{tabular}{@{}lrl@{}}
\toprule
Quantity & Value & Source \\
\midrule
$\srung$ ($\mu_c$)        & $0.01388338$ & ZC89 \\
$\sGam/\srung$ ($\mu_d/\mu_c$) & $36.9638$ & FIND-ZC89-01 (irrational) \\
$8\ln2/[\ln15\cdot\ln\rborn]$  & $36.9638$ & closed form, ZC89 \\
\bottomrule
\end{tabular}}
\end{center}

\textbf{Status:} Exact Lemma (inherits FIND-ZC89-01).
\end{formalbox}

%%-------------------------------------------------------------
\section{The Mind Layer Is a Q-Pair\quad(LEM-ZC90-02)}
\label{sec:lem02-zc90}
%%-------------------------------------------------------------

% [BRIDGE-PROSE: stakes before the box]
The mind layer is the harder instantiation, not because its structure is
unclear but because its magnitudes are uncalibrated. DEF-MB03 gives us
the \emph{kind} distinction in plain language and EQ-MB01 gives us the
signs and the inequality $\db\gg\dv$, but no paper has ever fixed the
\emph{number} $\db/\dv$. The lemma below therefore establishes the Q-pair
\emph{structure} on the mind layer while being honest that its
non-commensurability is asserted by the corpus (``of kind, not
magnitude'') rather than proved numerically --- a distinction we carry
forward openly as GAP-ZC90-01.

\needspace{8\baselineskip}
\begin{formalbox}[title={LEM-ZC90-02: The Mind-Layer Pair Is a Q-Pair
  ($R_0$, Near-Formal)}]

\textbf{Statement.} The pair
$(\,\mathcal{D}_{\rm mind};\,\text{visesa},\text{nibbedha})$ is a Q-pair,
where $\mathcal{D}_{\rm mind}$ is the Barami current (the forward
direction toward $\Sind$ in $D_{\rm KL}$), \emph{visesa} the continuous
generator, and \emph{nibbedha} the discontinuous generator.

\textbf{Derivation.}

\textit{Step 1 (conserved direction).}
By EQ-MB01, both \emph{visesa} ($-\dv$) and \emph{nibbedha} ($-\db$)
decrease $D_{\rm KL}(\tau_n)$ --- they move in the same forward sense
toward $\Sind$. DEF-MB04 defines Barami precisely as the accumulation
over this forward subset $\Cfourplus=\{\text{visesa},\text{nibbedha}\}$.
The Barami current is the conserved direction $\mathcal{D}_{\rm mind}$.

\textit{Step 2 (two generators of different kind).}
EQ-MB01: \emph{visesa} produces an ordinary forward increment $-\dv$;
\emph{nibbedha} produces a \emph{discontinuous} forward increment $-\db$
with $\db\gg\dv$. DEF-MB03 states the distinction is ``of kind, not of
magnitude'' --- i.e.\ \emph{nibbedha} is not a large \emph{visesa} but a
generator of a different type. This is clause~(2)--(3) of DEF-ZC90-01.

\textit{Step 3 (non-commensurability --- corpus-asserted).}
DEF-MB03 states $\mathcal{C}_4$ ``cannot be reduced to a continuous
variable.'' In Q-pair terms: \emph{nibbedha} is not realisable as a
finite iterate of \emph{visesa}, which is the weak form of
condition~\eqref{eq:qpair-zc90}. The \emph{numeric} ratio $\db/\dv$ that would
establish the strong (irrational) form is not fixed by any current paper;
this is the only clause not proved here, and it is recorded as
GAP-ZC90-01.

\hfill (Near-Formal: structure exact; numeric ratio open.)

\textbf{Status:} Near-Formal Lemma. Structure inherited verbatim from
DEF-MB03/04 + EQ-MB01; non-commensurability holds in weak form, strong
form open (GAP-ZC90-01).
\end{formalbox}

%%-------------------------------------------------------------
\section{The Cross-Layer Quantum-Pair Theorem\quad(THM-ZC90-01)}
\label{sec:thm-zc90}
%%-------------------------------------------------------------

% [BRIDGE-PROSE: the "so what", outsider-readable, UPG-30]
With both layers shown to be Q-pairs, the theorem is almost a tautology
--- and that is the point. The value of ZC90 is not a new computation but
a recognition: a feature that looked like a quirk of the force ladder
(``why is that ratio irrational rather than a clean integer?'') and a
feature that looked like a quirk of contemplative psychology (``why are
there two forward moves instead of one graded one?'') are the same
feature. The framework was stating one law twice without noticing. For an
outsider the upshot is this: CRF predicts that wherever it builds a
conserved direction, motion along it will come in exactly two
irreducible kinds --- a smooth quantum and a jump quantum --- and this is
now a stated, in-principle-falsifiable pattern rather than a coincidence.

\needspace{10\baselineskip}
\begin{formalbox}[title={THM-ZC90-01: Cross-Layer Quantum-Pair Theorem
  ($R_0$, Exact in structure)}]

\textbf{Statement.} The force-layer pair (LEM-ZC90-01) and the mind-layer
pair (LEM-ZC90-02) are instances of the \emph{same} structure
DEF-ZC90-01. Consequently the force-layer property ``$\sGam/\srung$ is
irrational'' and the mind-layer property ``\emph{visesa} and
\emph{nibbedha} differ in kind not magnitude'' are two readings of a
single structural statement:
\begin{equation}
  \boxed{\;
  \text{Q-pair}(\mathcal{D};g_c,g_d)\ \text{holds}\ \Longleftrightarrow\
  \frac{\mu_d}{\mu_c}\ \text{is non-commensurate}
  \;}
  \label{eq:thm-zc90}
\end{equation}
realised at the force layer by $(\srung,\sGam)$ with
$\mu_d/\mu_c=36.9638\notin\mathbb{Q}$, and at the mind layer by
$(\text{visesa},\text{nibbedha})$ with $\db/\dv$ non-commensurate
(strong form open, GAP-ZC90-01).

\textbf{Proof.}

\textit{Step 1 (both are Q-pairs).}
LEM-ZC90-01 and LEM-ZC90-02 verify every clause of DEF-ZC90-01 for the
force and mind layers respectively.

\textit{Step 2 (the defining property is the shared one).}
The single discriminating clause of DEF-ZC90-01 is the
non-commensurability condition~\eqref{eq:qpair-zc90}. At the force layer it is
proved (irrational ratio); at the mind layer it is the corpus statement
``of kind, not magnitude / cannot be reduced to a continuous variable.''
These are the same predicate applied on two layers.

\textit{Step 3 (no stronger identification claimed).}
The theorem identifies the \emph{structure}, not the
\emph{generators}: it does not assert $\srung=\dv$ or any numeric
equality across layers. Cross-layer numeric identity is neither needed
nor claimed (see Scope).

\hfill$\square$

\physmeaning{Forces and minds are built by CRF from the same primitive;
this theorem says they also \emph{move} by the same rule --- one smooth
step and one leap, never reducible to each other.}

\textbf{Status:} Exact in structure (the structural identification is
exact; the mind-layer numeric ratio remains open, GAP-ZC90-01).
\end{formalbox}

%%-------------------------------------------------------------
\section{The Gift: Barami as a Conserved Charge\quad(FIND-ZC90-01)}
\label{sec:find-zc90}
%%-------------------------------------------------------------

% [BRIDGE-PROSE: stakes before the box]
The theorem above came at the cost of a discarded hypothesis, and the
discard left a gift. Once we stop trying to read the four choices as a
group and instead read the two forward choices as quanta of a current,
DEF-MB04 stops being a mere definition and starts looking like a
conservation law: Barami is literally the accumulated count, weighted by
step size, of forward quanta emitted along the mind's history.

\needspace{8\baselineskip}
\begin{giftbox}[title={FIND-ZC90-01: Barami Is a Conserved Charge on the
  Forward 2-Set (Gift)}]

\textbf{Statement.} Barami (DEF-MB04) is the conserved charge of the
mind-layer Q-pair: the monotone accumulation along $\mathcal{D}_{\rm
mind}$ of quanta emitted by the two forward generators,
\[
  \mathrm{Barami}(\mathcal{M},\tau_N)
  =\big|\{n\le N: c_n\in\Cfourplus\}\big|\cdot
   \langle\dv,\db\rangle_{n\le N},
\]
exactly DEF-MB04, now \emph{interpreted}: the integer factor counts
quanta; the average-magnitude factor weights continuous and discontinuous
quanta of the one current.

\textbf{Why this is a gift (and why Klein-four failed).}
A Klein-four reading $\mathbb{Z}_2\times\mathbb{Z}_2$ treats all three order-2
subgroups symmetrically; nothing privileges any 2-element subset. But
DEF-MB04 \emph{does} privilege $\Cfourplus=\{\text{visesa},\text{nibbedha}\}$
--- it is the only subset on which Barami is defined. A symmetric group
cannot express a privileged direction; a conserved current can. The
asymmetry that breaks the group is exactly the orientation that defines
the charge. \emph{The refutation was the discovery.}

\physmeaning{What a mind accumulates over a life --- its Barami --- is the
total forward ``charge'' it has emitted, counting both its steady small
advances and its rare breakthroughs.}

\textbf{Status:} Near-Formal Finding (Gift-class --- a structural
reinterpretation of an existing definition; introduces no new number).
\end{giftbox}

%%-------------------------------------------------------------
\section{Conjecture and Open Gap\quad(CONJ/GAP-ZC90-01)}
\label{sec:conjgap-zc90}
%%-------------------------------------------------------------

% [BRIDGE-PROSE: stakes before the box]
The one thing the force layer has and the mind layer lacks is a number.
At the force layer the non-commensurability is not merely asserted but
computed: $36.9638=8\ln2/[\ln15\ln\rborn]$, every factor a $\Ztf$
quantity. It is natural to ask whether the mind layer hides an analogous
closed form. We state this as a conjecture, carefully labelled a
candidate, and register the missing number as the paper's open gap.

\begin{derivedbox}[title={CONJ-ZC90-01: The Mind Ratio Is a
  $\Ztf$-Structural Irrational (Candidate)}]
\textbf{Statement.} The mind-layer ratio $\db/\dv$ is, like the
force-layer $36.9638$, an irrational number expressible through
$\Ztf$-structural quantities ($\ln2$, $\ln15$, $d_s$, and a Born-chain
weight).

\textbf{Caution.} This is a \emph{candidate}, motivated only by
cross-layer analogy (THM-ZC90-01) and the shared $\Ztf$ substrate. No
derivation and no numeric estimate is offered. It must not be cited as
support for any prediction until GAP-ZC90-01 is closed.

\textbf{Status:} Candidate (conjecture).
\end{derivedbox}

\begin{openqbox}[title={GAP-ZC90-01: The Mind-Layer Step Ratio Is
  Undetermined}]
\textbf{Open item.} The numeric value of $\db/\dv$ (equivalently, the
strong-form non-commensurability of the mind-layer Q-pair) is not fixed
by any current CRF paper. EQ-MB01 fixes only $\db\gg\dv$ and the signs.

\textbf{What closing it would require.} A calibration of the agentive-mode
step magnitudes --- presumably from a Free-Energy / $D_{\rm KL}$ model of
the agentive trajectory (EQ-MB02, the $\Sigma$-Lagrange-constraint form)
--- yielding $\db/\dv$ in closed form, testing CONJ-ZC90-01.

\textbf{Target:} ZC91.
\end{openqbox}

%%-------------------------------------------------------------
\section{PM-ZC90-01: The Klein-Four Reading (Productive Scar)}
\label{sec:pm-zc90}
%%-------------------------------------------------------------

\begin{derivedbox}[title={PM-ZC90-01: $\mathcal{C}_4$ as Klein-Four
  $\mathbb{Z}_2\times\mathbb{Z}_2$ (Failure Stance)}]
\textbf{What happened.} A probe assigned the four choices to
$\mathbb{Z}_2\times\mathbb{Z}_2$ via two binary axes (direction; continuous vs
discontinuous), and reported a clean bijection onto the four cosets.

\textbf{Why it was not an error.} Phase Misalignment: the probe forced
\emph{ṭhiti} into a ``backward'' cell, but the corpus (DEF-MB03) makes
\emph{ṭhiti} a passive \emph{anicca} drift ($+\epsilon_t$), not a directed
choice --- so the second axis is 3-vs-1, not 2-vs-2, and Klein-four
cannot hold. The decisive fact is DEF-MB04: Barami privileges the forward
pair, which a symmetric Klein group cannot represent.

\textbf{Preservation.} Kept as a productive scar: refuting the group
reading is what exposed the conserved-current reading (FIND-ZC90-01) and
the cross-layer theorem (THM-ZC90-01).

\textbf{Pattern type:} tautology test (the probe tested a proposed
algebraic definition, not a physical claim).

\textbf{Lesson --- PWA-10:} when a probe proposes an algebraic structure,
audit it against the corpus's \emph{distinguished} objects (here Barami)
before accepting the symmetry; a privileged subset refutes a symmetric
group.
\end{derivedbox}

%%=============================================================
%% ATOMIC REGISTRY (MANDATORY)
%%=============================================================
\begin{formalbox}[title={Atomic Registry --- ZC90}]
\begin{center}
\renewcommand{\arraystretch}{1.3}
\small
\begin{tabular}{lp{8.0cm}l}
\toprule
\textbf{ID} & \textbf{Description} & \textbf{Status} \\
\midrule
DEF-ZC90-01
  & Quantum-Pair: conserved direction + continuous + discontinuous
    generator, non-commensurate
  & Definition \\[3pt]
LEM-ZC90-01
  & Force-layer pair $(\srung,\sGam)$ is a Q-pair (irrational ratio)
  & Exact \\[3pt]
LEM-ZC90-02
  & Mind-layer pair (visesa, nibbedha) is a Q-pair on the Barami current
  & Near-Formal \\[3pt]
THM-ZC90-01
  & Cross-Layer Quantum-Pair Theorem: both layers, one signature
  & Exact (structure) \\[3pt]
FIND-ZC90-01
  & Barami is the conserved charge of the mind-layer Q-pair
  & Near-Formal (Gift) \\[3pt]
CONJ-ZC90-01
  & $\db/\dv$ is a $\Ztf$-structural irrational
  & Candidate \\[3pt]
GAP-ZC90-01
  & Numeric value of $\db/\dv$ undetermined; target ZC91
  & Open \\[3pt]
PM-ZC90-01
  & $\mathcal{C}_4$-as-Klein-four refuted (pattern: tautology test)
  & Scar \\[3pt]
\midrule
\multicolumn{3}{l}{\emph{Gap-status changes from prior papers:} none ---
  no prior gap closed (structural bridge paper).} \\
\bottomrule
\end{tabular}
\end{center}
\end{formalbox}

%%  REFRAME MARKER (used — a genuine proof-direction change occurred)
%% UPG-23: the paper's key move was "we thought Klein-four → it turned out
%% conserved current", which changed the proof route. Qualifies (cf. ZC71).
\begin{derivedbox}[title={Reframe: $\mathcal{C}_4$ as a group
  $\to$ $\mathcal{C}_4^{+}$ as a conserved current}]
\textbf{Before this paper:} We expected the four-choice $\mathcal{C}_4$ to
carry a finite group structure (Klein-four candidate).

\textbf{After this paper:} The correct framing is a \emph{pointed
conserved current} on the forward pair $\Cfourplus$; the two forward
choices are quanta, not cosets.

\textbf{Proof route opened:} the Q-pair definition and the cross-layer
identification with the force layer (THM-ZC90-01).

\textbf{Proof route closed:} the dead end of seeking a group
homomorphism / Cayley structure for $\mathcal{C}_4$ (PM-ZC90-01).
\end{derivedbox}

%%=============================================================
\section{R-Layer Research Note}
\label{sec:rlayer-zc90}
%%=============================================================
\begin{layerbox}[title={R-Layer Assignments for ZC90}]

\textbf{Layer assignments:}

{\small\renewcommand{\arraystretch}{1.25}
\begin{center}
\begin{tabular}{@{}llll@{}}
\toprule
Atomic unit & R-layer & Cost $T$ & Source \\
\midrule
DEF-ZC90-01 & $R_0$ & $0$ & structural definition \\
LEM-ZC90-01 & $R_0$ & $0$ & algebraic from FIND-ZC89-01 \\
LEM-ZC90-02 & $R_0$ & $0$ & structural from DEF-MB03/04 \\
THM-ZC90-01 & $R_0$ & $0$ & structural identification \\
GAP-ZC90-01 & $R_0$ & --- & magnitude calibration (not a bridge cost) \\
\bottomrule
\end{tabular}
\end{center}}

\textbf{Two-layer floor (PWA-8):}
\begin{itemize}[leftmargin=2em, noitemsep]
  \item Spectral layer ($R_0$): $\epscan=0.007550$ --- not used here.
  \item Bridge layer ($R_0{\to}R_{\max}$): $\epsuniv=0.007599$ --- enters
    only \emph{inside} the inherited ZC89 ratio (via $\epsEM,\epsW$ in
    $\srung$); ZC90 introduces no new floor. Separation $0.66\%$
    irreducible (THM-ZC71-01).
\end{itemize}

\textbf{Key R-layer remark.} All ZC90 results are at $R_0$ and purely
structural; no Born-chain bridge cost is computed in this paper. The only
R-layer content is inherited intact from ZC89.
\end{layerbox}

%%=============================================================
\section{Genealogy Map}
\label{sec:genealogy-zc90}
%%=============================================================
\begin{genealogybox}[title={How ZC90's Key Objects Trace Back}]

\textbf{Branch A --- Force-layer step structure:}
\begin{itemize}[leftmargin=2em, noitemsep]
  \item \textbf{Part~B / TLS floor $\epscan$}: the instability floor.
  \item $\to$ \textbf{ZC73 / FIND-ZC73-01}: the Gamma ladder, step
    $\ln2/\ln15$.
  \item $\to$ \textbf{ZC88 / THM-ZC88-01}: floor lattice and Gamma ladder
    are one co-AP under the log-map, step $\srung$.
  \item $\to$ \textbf{ZC89 / FIND-ZC89-01}: the two steps' ratio
    $36.9638$ is irrational --- a continuous/discontinuous pair.
  \item $\to$ \textbf{ZC90 / LEM-ZC90-01}: this pair is a Q-pair.
\end{itemize}

\textbf{Branch B --- Mind-layer choice structure:}
\begin{itemize}[leftmargin=2em, noitemsep]
  \item \textbf{Part~F / DEF-MB02}: $\delta$ in Agentive Mode.
  \item $\to$ \textbf{Part~F / DEF-MB03}: the four-way choice
    $\mathcal{C}_4$; \emph{visesa}/\emph{nibbedha} ``of kind not magnitude.''
  \item $\to$ \textbf{Part~F / DEF-MB04}: Barami defined on the forward
    pair $\Cfourplus$.
  \item $\to$ \textbf{ZC90 / LEM-ZC90-02, FIND-ZC90-01}: the forward pair
    is a Q-pair; Barami is its conserved charge.
\end{itemize}

\textbf{Convergence at ZC90:}
Branch~A reached a continuous/discontinuous non-commensurate pair by
\emph{proof} (irrational ratio); Branch~B carried the same pair by
\emph{definition} (``kind not magnitude'') since Part~F. DEF-ZC90-01
names the common structure and THM-ZC90-01 identifies the two as one.

ZC90's contribution: it shows that a force-layer theorem (ZC89) and a
mind-layer definition (DEF-MB03) are the same structural statement, and
extracts Barami-as-conserved-charge (FIND-ZC90-01) from the refutation of
a group reading --- a unification no prior paper made.
\end{genealogybox}

%%=============================================================
\section{Closing Sentence}
%%=============================================================
%% UPG-22 — TONE B (Contemplative): this is a structural bridge paper;
%% it identifies a shared structure and opens GAP-ZC90-01. Tone B selected
%% (declared in header). Precedents: ZC72, ZC75.

\bigskip
\begin{center}
\textit{%
  We came hunting a group and found a current: two forward movements,
  one a steady advance, one a sudden breaking-through, sharing a single
  direction yet measuring each other by no common rule --- the same law
  the forces obey in the irrational $36.9638$.\\[0.3em]
  The structure is one; the mind's number belongs to ZC91.
}
\end{center}

% Governance Note: This document follows Upgrade Protocol v1.3
% and CRF Style Guide v3.7.


% [BRIDGE-PROSE: integration-added, Style Guide v3.7 §31.6 / UPG-30]
\noindent
ZC90 pairs the two mind-steps but leaves open (GAP-ZC90-01) whether they are
truly distinct kinds or just two sizes of one kind --- a question with a
conjecture (CONJ-90-01) attached. If they collapsed into one kind, the whole
mind-layer structure would reduce to a rescaling of the force layer. ZC91
settles it: the two steps are irrational-of-kind, retiring the conjecture and
showing the mind layer is structurally its own thing.
\section{Background: The Question That Dissolved}
\label{sec:background-zc91}
%%=============================================================

%% ── [v3.6 §31.2 / UPG-25] READER'S ON-RAMP ───────
CRF describes a conscious agent as choosing, at each moment, among four
mental movements; two of them move ``forward'' toward a stable inner
ground. One is an ordinary advance (\emph{visesa}), the other a sudden
breakthrough (\emph{nibbedha}). A previous paper (ZC90) showed these two
play the same structural role that two interaction steps play in the
physics sector, and asked the obvious next question: how much bigger is
the breakthrough than the ordinary step --- what is the number
$\db/\dv$? This paper answers that the number does not exist as a fixed
constant, and that this is not a failure but the point: the framework's
own definition of the breakthrough \emph{forbids} it from being a clean
multiple of the ordinary step. ``A different kind, not a bigger size'' is
exactly the statement that no fixed ratio can capture the difference. The
sections below turn that plain observation into a proof.

%% ── [NEW BLOCK 1] NUMERICAL SUMMARY ────────────────────────
\begin{keybox}[title={Numerical Summary --- ZC91 (T-canonical)}]
{\footnotesize\renewcommand{\arraystretch}{1.3}
\begin{center}
\begin{tabular}{@{}p{5.4cm}p{3.2cm}p{3.6cm}@{}}
\toprule
\textbf{Quantity} & \textbf{Value} & \textbf{Source / Status} \\
\midrule
$\Nstar=\db/\dv$ (mind ratio)
  & irrational; \emph{no constant}
  & \textbf{THM/COR-ZC91 (new)} \\
force-layer analogue $\sGam/\srung$
  & $36.9638\ldots$ (irrational)
  & ZC89 (FIND-ZC89-01) \\
$\lim_{\Barami\to\infty}(\dv,\db)$
  & $\to 0$ (both)
  & EQ-MB03 (Part~F) \\
ratio limit $\lim\Nstar$
  & open
  & GAP-ZC91-01 (new) \\
\bottomrule
\end{tabular}
\end{center}}
\end{keybox}

%% ── PHYSICS / FRAMEWORK CONTEXT ─────────────────────────────
\subsection{What ZC90 Left Open, and What the Corpus Actually Fixes}

ZC90 established that the forward pair $(\text{visesa},\text{nibbedha})$
is a Quantum-Pair (DEF-ZC90-01): a conserved direction carrying a
continuous and a discontinuous generator. It proved the strong
non-commensurability only at the force layer (the irrational
$\sGam/\srung=36.9638$, FIND-ZC89-01) and left the mind-layer strong form
open as GAP-ZC90-01, conjecturing (CONJ-ZC90-01) that $\Nstar$ might be a
specific $\Ztf$-structural irrational, by analogy.

Tracing upstream into the Mind/Barami corpus (the rule of ZC81: go to
where a quantity is \emph{defined}, not where it is used) shows what is
and is not pinned. EQ-MB01 fixes the signs and the ordering $\db\gg\dv$.
EQ-MB03 fixes the asymptotic limit: both $\dv$ and $\db$ tend to zero as
Barami grows without bound, recovering gradient flow (P-MB04). And P-MB02
states something decisive: $\dv$ at a given moment is a \emph{function of
accumulated Barami}, so two minds in identical present states have
different step magnitudes. The corpus therefore fixes ordering and limit
but supplies no numerical anchor, and --- crucially --- makes the ratio a
property of a mind's history, not a universal constant. The object
CONJ-ZC90-01 went looking for may not exist.

\subsection{What Was Already Known Before ZC91}

\begin{keybox}[title={Pre-ZC91 Status}]
\textbf{Known:}
\begin{itemize}[leftmargin=2em, noitemsep]
  \item Force layer: $\sGam/\srung=36.9638$ is exact and irrational
    (FIND-ZC89-01).
  \item Mind layer: $(\text{visesa},\text{nibbedha})$ is a Q-pair; strong
    non-commensurability open (GAP-ZC90-01).
  \item Mind layer: \emph{visesa}/\emph{nibbedha} ``of kind, not
    magnitude''; $\mathcal{C}_4$ irreducible to a continuous variable
    (DEF-MB03).
  \item Mind layer: $\db\gg\dv>0$ (EQ-MB01); both $\to0$ as
    $\Barami\to\infty$ (EQ-MB03); $\dv$ is Barami-history-dependent
    (P-MB02).
\end{itemize}

\textbf{Open:}
\begin{itemize}[leftmargin=2em, noitemsep]
  \item Strong-form non-commensurability of the mind-layer Q-pair
    (GAP-ZC90-01).
  \item Whether $\Nstar=\db/\dv$ has a $\Ztf$-structural closed form
    (CONJ-ZC90-01).
\end{itemize}

ZC91 closes GAP-ZC90-01 (THM-ZC91-01) and resolves CONJ-ZC90-01
negatively (COR-ZC91-02). What it cannot close --- the Barami-history
functional --- it re-opens precisely as GAP-ZC91-01.
\end{keybox}

%% ── [NEW BLOCK 2] SCOPE EXCLUSION ───────────────────────────
\begin{keybox}[title={What ZC91 Does NOT Do}]
\begin{itemize}[leftmargin=2em, noitemsep]
  \item Does not compute a numeric value for $\Nstar$ (no constant value
    exists; COR-ZC91-02).
  \item Does not modify DEF-MB01--07 or EQ-MB01--04 (used verbatim).
  \item Does not modify ZC89/ZC90 or any force-layer result.
  \item Does not derive the Barami-history functional (GAP-ZC91-01).
  \item Does not claim the mind-layer and force-layer irrationals are
    equal or related in value --- only that both are \emph{forced}
    irrational.
\end{itemize}
\end{keybox}

%% ── [NEW BLOCK 3] COMPANION PAPERS ─────────────────────────
\begin{keybox}[title={Companion Papers --- ZC91}]
{\footnotesize\renewcommand{\arraystretch}{1.3}
\begin{center}
\begin{tabular}{@{}p{6.1cm}p{6.4cm}@{}}
\toprule
\textbf{Paper} & \textbf{What it provides to ZC91} \\
\midrule
\texttt{CRF\_ZC90\_CrossLayerQuantumPair\_v1.tex}
  & DEF-ZC90-01 (Q-pair); GAP-ZC90-01, CONJ-ZC90-01 ---
    \textbf{the items this paper resolves} \\
\texttt{CRF\_ZC89\_LayerSpecificity\_v3.tex}
  & FIND-ZC89-01: the force-layer irrational, the independent analogue \\
\texttt{CRF\_Complete\_v17\_16\_Part\_F.tex}
  & DEF-MB03, EQ-MB01/02/03, P-MB02/04 --- \textbf{KEY PILLARS} \\
\texttt{CRF\_Complete\_Index\_v17\_16.tex}
  & Part~F (Mind/Barami) context \\
\bottomrule
\end{tabular}
\end{center}}
\end{keybox}

%% ── DISCOVERY NOTE (promoted) ──
\begin{derivedbox}[title={Discovery Note}]
This paper was commissioned to \emph{find} the closed form of $\Nstar$. A
bidirectional probe --- forward from the MB axioms, backward from the
ZC89 force-layer analogue --- found the reverse. Forward: the corpus
pins only ordering and limit, and P-MB02 makes the ratio
history-dependent. Backward: trying to write $\db=\Nstar\dv$ with
$\Nstar$ a clean number \emph{is} a magnitude statement, which DEF-MB03
forbids. The closed form we sought could not exist, and the reason it
could not exist is the theorem. The proof direction inverted mid-session;
see the Reframe Marker.
\end{derivedbox}

%% ── PRE-WORK AUDIT ──────────────────────────────────────────
\subsection{Pre-Work Audit (PWA-1 through PWA-10)}

\begin{formalbox}[title={Pre-Work Audit Record --- ZC91}]
Scanned before writing (Upgrade Protocol v1.3, PWA-1 through PWA-10):
\begin{itemize}[leftmargin=2em, noitemsep]
  \item ZC90/DEF-ZC90-01, GAP-ZC90-01, CONJ-ZC90-01 --- the items resolved.
  \item Part~F/DEF-MB03 --- \textbf{KEY PILLAR} (the kind/magnitude clause).
  \item Part~F/EQ-MB01, EQ-MB03, P-MB02, P-MB04 --- the upstream constraints,
    read verbatim (not reconstructed).
  \item ZC89/FIND-ZC89-01 --- the independent force-layer analogue.
\end{itemize}

\textbf{PWA-DOWNSTREAM} (ZC81 rule): traced $\dv,\db$ to their definition
(EQ-MB01) and their dynamics (EQ-MB03, P-MB02) \emph{before} proposing any
value --- which is what revealed that a constant is the wrong object. No
value was forced onto a history-dependent quantity.\\
\textbf{PWA-6} (macro convention): paper macros all
\texttt{\textbackslash providecommand}; no Preamble collision.\\
\textbf{PWA-7/8/9}: not triggered (no R-layer bridge, no new floor, no
spectral convention).\\
\textbf{PWA-10} (tautology audit): THM-ZC91-01 reads a consequence
\emph{out of} a definition (DEF-MB03). This is flagged and addressed
head-on in the theorem's Remark --- it is an equivalence, made explicit,
not a hidden tautology dressed as physics.\\
\textbf{No duplicate atomic units found.}
\end{formalbox}

\newpage

%%=============================================================
%% ATOMIC UNITS — THM → COR → FIND → GAP → OBS
%%=============================================================

%%-------------------------------------------------------------
\section{The Irrationality-of-Kind Theorem\quad(THM-ZC91-01)}
\label{sec:thm-zc91}
%%-------------------------------------------------------------

% [BRIDGE-PROSE: stakes before the box, UPG-30]
The whole paper turns on a single move, and it is worth stating in plain
words before any symbols. To say ``the breakthrough is the same kind of
thing as an ordinary step, only $\Nstar$ times bigger'' is to say the two
differ in magnitude. DEF-MB03 explicitly denies this: the difference is
of kind. The only way a quantitative ratio can encode a difference that
is \emph{not} of magnitude is for that ratio to be unreachable by any
finite counting of the smaller unit --- that is, irrational. The theorem
is the careful version of this sentence, and its one delicate point,
which we do not hide, is that it draws a quantitative conclusion out of a
qualitative definition. We argue that this is legitimate precisely
because DEF-MB03's clause is itself a claim about quantity (it denies
one).

\needspace{6\baselineskip}
\begin{formalbox}[title={THM-ZC91-01: Irrationality-of-Kind ($R_0$, Exact)}]

\textbf{Statement.} Let $\dv,\db>0$ be the \emph{visesa} and
\emph{nibbedha} step magnitudes of EQ-MB01 at any fixed moment, and
$\Nstar=\db/\dv$. Then DEF-MB03's clause --- that \emph{visesa} and
\emph{nibbedha} differ ``of kind, not of magnitude'', equivalently that
$\mathcal{C}_4$ ``cannot be reduced to a continuous variable'' --- is
equivalent to
\begin{equation}
  \boxed{\;\Nstar=\frac{\db}{\dv}\notin\mathbb{Q}.\;}
  \label{eq:thm-zc91}
\end{equation}

\textbf{Proof.}

\textit{Step 1 ($\Nstar\in\mathbb{Q}\Rightarrow$ magnitude difference).}
Suppose $\Nstar=p/q$ with $p,q\in\mathbb{N}$ coprime. Then
$q\,\db=p\,\dv$: a whole number $q$ of \emph{nibbedha} increments equals a
whole number $p$ of \emph{visesa} increments along the same conserved
direction $\mathcal{D}_{\rm mind}$ (LEM-ZC90-02). A breakthrough is then
exactly $p/q$ ordinary steps; \emph{nibbedha} is reducible to a counted
multiple of \emph{visesa}. This is a difference of \emph{magnitude} only
--- precisely what DEF-MB03 denies.

\textit{Step 2 ($\Nstar\in\mathbb{Q}\Rightarrow$ continuous-variable
reduction).} If $q\,\db=p\,\dv$, the two generators lie on one rational
lattice $\{k\,\dv/q : k\in\mathbb{Z}\}$, a single graded one-parameter
family. The four-way choice then \emph{is} reducible to position on a
continuous (here, finely discretised) variable --- contradicting the
explicit DEF-MB03 clause that it is not.

\textit{Step 3 (contrapositive and converse).} By Steps 1--2, the
DEF-MB03 clause implies $\Nstar\notin\mathbb{Q}$. Conversely, if
$\Nstar\notin\mathbb{Q}$, no finite $p,q$ satisfy $q\,\db=p\,\dv$, so
\emph{nibbedha} is not any finite multiple of \emph{visesa}: the two are
incommensurable and the distinction is necessarily one of kind, not
magnitude. The equivalence~\eqref{eq:thm-zc91} holds.

\hfill$\square$

\textbf{Remark (the PWA-10 point, addressed openly).} This theorem
extracts a numerical property (irrationality) from a definitional clause
(``kind, not magnitude''). That is admissible here because the clause is
\emph{already} a quantitative assertion --- it denies the existence of a
finite ratio. We are not smuggling physics out of a tautology; we are
making explicit the arithmetic content that DEF-MB03 stated in words. The
force layer reaches the same conclusion by entirely different means
(ZC89's ladder arithmetic), which is the independent check (FIND-ZC91-01).

\physmeaning{A breakthrough is not a big ordinary step. No matter how many
ordinary steps you take, you never exactly equal one breakthrough --- and
that ``never exactly'' is what irrational means here.}

\textbf{Status:} Exact Theorem.
\end{formalbox}

%%-------------------------------------------------------------
\section{Corollaries\quad(COR-ZC91-01, COR-ZC91-02)}
\label{sec:cor-zc91}
%%-------------------------------------------------------------

% [BRIDGE-PROSE]
Two consequences follow immediately, one positive and one negative. The
positive one closes the gap ZC90 left. The negative one retires the
conjecture ZC90 raised --- not as a mistake, but as a question now
answered in the direction nobody bet on.

\needspace{8\baselineskip}
\begin{formalbox}[title={COR-ZC91-01: GAP-ZC90-01 Closed ($R_0$, Exact)}]
\textbf{GAP-ZC90-01} (ZC90): the strong-form non-commensurability of the
mind-layer Q-pair ($\Nstar\notin\mathbb{Q}$) was open, pending a numeric
calibration.

\textbf{Resolution.} THM-ZC91-01 proves $\Nstar\notin\mathbb{Q}$ directly
from DEF-MB03, with no calibration required. The strong form holds.

\textbf{Axiom chain:}
\[
  \delta \xrightarrow{\text{Agentive Mode}} \mathcal{C}_4
  \xrightarrow{\text{DEF-MB03 (kind}\neq\text{magnitude)}}
  \Nstar\notin\mathbb{Q}.
\]
\textbf{Status:} GAP-ZC90-01 \textbf{CLOSED}. $\checkmark$
\end{formalbox}

\needspace{8\baselineskip}
\begin{formalbox}[title={COR-ZC91-02: CONJ-ZC90-01 Resolved Negative
  ($R_0$, Exact)}]
\textbf{CONJ-ZC90-01} (ZC90, Candidate): that $\Nstar$ is a specific
$\Ztf$-structural irrational constant (by analogy to the force-layer
$36.9638$).

\textbf{Resolution.} Negative. P-MB02 states $\dv$ (hence $\Nstar$) is a
function of accumulated Barami: two minds in identical present states have
different $\dv$. A single universal \emph{constant} $\Nstar$ therefore
does not exist; the analogy to the force-layer constant fails at the
level of object type (the force ratio is a fixed structural number; the
mind ratio is a history functional). THM-ZC91-01's irrationality holds
\emph{pointwise} (at each moment, for each mind) but there is no constant
to assign a closed form to.

\textbf{Status:} CONJ-ZC90-01 \textbf{RESOLVED NEGATIVE}. Retired
conjecture (correctly declared Candidate in ZC90; not a scar). Provenance:
OBS-ZC91-01.
\end{formalbox}

%%-------------------------------------------------------------
\section{The Gift: Two Independent Routes to One Necessity\quad(FIND-ZC91-01)}
\label{sec:find-zc91}
%%-------------------------------------------------------------

% [BRIDGE-PROSE: the "so what", outsider-readable, UPG-30]
Here is the part worth pausing on. The force layer was forced to an
irrational ratio by arithmetic: two ladders whose steps, built from
$\ln2$, $\ln15$, and a Born-chain weight, simply do not divide one
another (ZC89). The mind layer is forced to an irrational ratio by
meaning: a definition that insists a breakthrough is a different kind of
event, not a larger one (DEF-MB03). These two arguments share no
machinery --- one is number theory over the group order, the other is a
semantic clause about consciousness. That two completely unrelated routes
arrive at the same necessity is the kind of redundancy that, in physics,
usually signals a real law rather than an accident. CRF appears to be
saying: \emph{wherever it builds a conserved direction with a smooth
quantum and a jump quantum, the two are incommensurable} --- and it says
so twice, in two languages, without having arranged to.

\needspace{8\baselineskip}
\begin{giftbox}[title={FIND-ZC91-01: Independent-Route Redundancy (Gift)}]
\textbf{Statement.} The non-commensurability of a CRF Quantum-Pair is
over-determined: it is forced
\begin{itemize}[leftmargin=2em, noitemsep]
  \item at the \textbf{force layer} by ladder arithmetic ---
    $\sGam/\srung=8\ln2/[\ln15\ln\rborn]$ is irrational (ZC89), a
    statement about the group order $\Ztf$ and the Born chain;
  \item at the \textbf{mind layer} by the kind/magnitude distinction ---
    $\Nstar\notin\mathbb{Q}$ (THM-ZC91-01), a statement about the meaning
    of \emph{nibbedha} in DEF-MB03.
\end{itemize}
The two derivations share no premises. Their agreement promotes ``a CRF
Q-pair has an irrational ratio'' from a per-layer observation to a
candidate structural law of the framework.

\textbf{Why this is a gift.} We set out to compute a number and instead
found a redundancy: the same necessity proved twice by disjoint means.
Redundant proof of the same fact, from independent axioms, is evidence
the fact is structural rather than coincidental.

\physmeaning{The universe of CRF seems to dislike clean ratios between a
smooth process and a sudden one --- and it enforces that dislike once with
arithmetic and once with meaning, never having coordinated the two.}

\textbf{Status:} Near-Formal Finding (Gift-class). Elevation to a stated
structural law awaits a third independent instance (a research direction,
not a claim).
\end{giftbox}

%%-------------------------------------------------------------
\section{Open Gap and Resolution Note\quad(GAP-ZC91-01, OBS-ZC91-01)}
\label{sec:gapobs-zc91}
%%-------------------------------------------------------------

% [BRIDGE-PROSE]
Closing the strong form sharpens, rather than ends, the inquiry. We now
know the ratio is irrational at every moment but is not a constant; what
governs how it varies is the genuinely open question, and it is a more
precise question than the one ZC90 asked.

\begin{openqbox}[title={GAP-ZC91-01: The Barami-History Functional}]
\textbf{Open item.} EQ-MB03 and P-MB02 make $\dv$ and $\db$ functions of
accumulated Barami, both tending to zero as $\Barami\to\infty$. Open:
\begin{itemize}[leftmargin=2em, noitemsep]
  \item the functional form $\dv(\Barami),\db(\Barami)$;
  \item whether $\Nstar(\Barami)=\db/\dv$ has a finite limit as
    $\Barami\to\infty$ (the two magnitudes vanish; their ratio need not),
    and whether that limit --- if it exists --- is structural.
\end{itemize}

\textbf{What closing it would require.} A Free-Energy / $D_{\rm KL}$
calibration of the agentive trajectory via EQ-MB02 (the
$\Sigma$-Lagrange-constraint Lagrangian), yielding the $\Barami$-dependence
of the step magnitudes. This is the first mind-layer result that would
need a dynamical model, not only the static definitions.

\textbf{Target:} ZC92.
\end{openqbox}

\begin{derivedbox}[title={OBS-ZC91-01: Resolution Note for CONJ-ZC90-01}]
For provenance: CONJ-ZC90-01 was introduced in ZC90 explicitly as a
Candidate, with the caution that it must not support any prediction until
tested. ZC91 tests it and resolves it negatively (COR-ZC91-02). This is
the conjecture lifecycle working as designed --- a labelled candidate,
raised and then settled --- and is recorded as an observation, not a
Phase-Misalignment scar.
\end{derivedbox}

%%=============================================================
%% ATOMIC REGISTRY (MANDATORY)
%%=============================================================
\begin{formalbox}[title={Atomic Registry --- ZC91}]
\begin{center}
\renewcommand{\arraystretch}{1.3}
\small
\begin{tabular}{lp{8.0cm}l}
\toprule
\textbf{ID} & \textbf{Description} & \textbf{Status} \\
\midrule
THM-ZC91-01
  & Irrationality-of-Kind: ``kind not magnitude'' $\Leftrightarrow$
    $\db/\dv\notin\mathbb{Q}$
  & Exact \\[3pt]
COR-ZC91-01
  & GAP-ZC90-01 closed: mind-layer strong non-commensurability proved
  & Closed $\checkmark$ \\[3pt]
COR-ZC91-02
  & CONJ-ZC90-01 resolved negative: no constant closed form ($\Nstar$ is
    a history functional)
  & Closed $\checkmark$ \\[3pt]
FIND-ZC91-01
  & Independent-route redundancy: force (arithmetic) + mind (semantic)
    both force irrationality
  & Near-Formal (Gift) \\[3pt]
GAP-ZC91-01
  & Barami-history functional $\dv(\Barami),\db(\Barami)$; ratio limit;
    target ZC92
  & Open \\[3pt]
OBS-ZC91-01
  & Resolution note for CONJ-ZC90-01 (lifecycle, not a scar)
  & Observation \\[3pt]
\midrule
\multicolumn{3}{l}{\emph{Gap-status changes from prior papers:}} \\
GAP-ZC90-01 & mind-layer strong non-commensurability; closed via
  COR-ZC91-01 & Closed $\checkmark$ \\
CONJ-ZC90-01 & ratio as $\Ztf$ constant; resolved negative via COR-ZC91-02
  & Resolved $-$ \\
\bottomrule
\end{tabular}
\end{center}
\end{formalbox}

%%  REFRAME MARKER (used — genuine proof-direction inversion; cf. ZC34/ZC71)
\begin{derivedbox}[title={Reframe: ``find the closed form''
  $\to$ ``prove no closed form can exist''}]
\textbf{Before this paper:} The task (inherited from GAP-ZC90-01) was to
\emph{compute} $\Nstar=\db/\dv$, expecting a $\Ztf$-structural value.

\textbf{After this paper:} The correct framing is that no constant value
exists (COR-ZC91-02), and the strong-form non-commensurability is
\emph{proved from the definition} (THM-ZC91-01) rather than calibrated.

\textbf{Proof route opened:} reading DEF-MB03's ``kind not magnitude'' as
the arithmetic statement $\Nstar\notin\mathbb{Q}$; the
independent-route gift (FIND-ZC91-01).

\textbf{Proof route closed:} the dead end of fitting a single constant to
a Barami-history-dependent quantity (the ZC81-scar pattern, avoided here
by going upstream first).
\end{derivedbox}

%%=============================================================
\section{R-Layer Research Note}
\label{sec:rlayer-zc91}
%%=============================================================
\begin{layerbox}[title={R-Layer Assignments for ZC91}]

\textbf{Layer assignments:}

{\small\renewcommand{\arraystretch}{1.25}
\begin{center}
\begin{tabular}{@{}llll@{}}
\toprule
Atomic unit & R-layer & Cost $T$ & Source \\
\midrule
THM-ZC91-01 & $R_0$ & $0$ & definitional/arithmetic (DEF-MB03) \\
COR-ZC91-01 & $R_0$ & $0$ & follows THM-ZC91-01 \\
COR-ZC91-02 & $R_0$ & $0$ & follows P-MB02 \\
FIND-ZC91-01 & $R_0$ & $0$ & cross-layer comparison \\
GAP-ZC91-01 & $R_0$ & --- & dynamical calibration (not a bridge cost) \\
\bottomrule
\end{tabular}
\end{center}}

\textbf{Two-layer floor (PWA-8):}
\begin{itemize}[leftmargin=2em, noitemsep]
  \item Spectral layer ($R_0$): $\epscan=0.007550$ --- not used.
  \item Bridge layer: $\epsuniv=0.007599$ --- enters only inside the
    inherited ZC89 ratio cited in FIND-ZC91-01; ZC91 introduces no floor.
    Separation $0.66\%$ irreducible (THM-ZC71-01).
\end{itemize}

\textbf{Key R-layer remark.} All results are at $R_0$ and structural; the
only R-layer content is the ZC89 ratio cited as the independent analogue,
inherited intact.
\end{layerbox}

%%=============================================================
\section{Genealogy Map}
\label{sec:genealogy-zc91}
%%=============================================================
\begin{genealogybox}[title={How ZC91's Key Objects Trace Back}]

\textbf{Branch A --- The mind-layer definition that became a theorem:}
\begin{itemize}[leftmargin=2em, noitemsep]
  \item \textbf{Part~F / DEF-MB03}: \emph{visesa}/\emph{nibbedha} ``of
    kind, not magnitude.''
  \item $\to$ \textbf{ZC90 / LEM-ZC90-02}: the forward pair is a Q-pair;
    strong non-commensurability left open (GAP-ZC90-01).
  \item $\to$ \textbf{ZC91 / THM-ZC91-01}: the clause \emph{is} the
    statement $\Nstar\notin\mathbb{Q}$; GAP-ZC90-01 closed.
\end{itemize}

\textbf{Branch B --- The force-layer analogue (independent check):}
\begin{itemize}[leftmargin=2em, noitemsep]
  \item \textbf{ZC88 / THM-ZC88-01}: floor and Gamma ladders are one
    co-AP.
  \item $\to$ \textbf{ZC89 / FIND-ZC89-01}: their step ratio
    $36.9638$ is irrational, by group-order arithmetic.
  \item $\to$ \textbf{ZC91 / FIND-ZC91-01}: this irrationality and the
    mind-layer's are forced by disjoint logic --- a redundancy.
\end{itemize}

\textbf{Branch C --- Why no constant exists:}
\begin{itemize}[leftmargin=2em, noitemsep]
  \item \textbf{Part~F / P-MB02}: $\dv$ is a function of accumulated
    Barami.
  \item \textbf{Part~F / EQ-MB03}: $\dv,\db\to0$ as $\Barami\to\infty$.
  \item $\to$ \textbf{ZC91 / COR-ZC91-02}: $\Nstar$ is a history
    functional, not a constant; CONJ-ZC90-01 negative.
\end{itemize}

\textbf{Convergence at ZC91:}
Branch~A turns a semantic clause into an arithmetic theorem; Branch~B
supplies the independent confirmation; Branch~C explains why the number
the task asked for is not a number at all but a functional.

ZC91's contribution: it closes a non-commensurability gap by proving the
result \emph{out of a definition}, shows the sought constant cannot exist,
and exposes a two-route redundancy that no single-layer paper could see.
\end{genealogybox}

%%=============================================================
\section{Closing Sentence}
%%=============================================================
%% UPG-22 — TONE A (Triumphant: a gap is fully closed, a theorem proved),
%% with one forward-pointing line (the gap re-opens sharper). Declared in
%% header. Precedent for "closure + sharpened successor": ZC69/ZC70 style.

\bigskip
\begin{center}
\textit{%
  We went to weigh the breakthrough against the step, and found that no
  scale exists that could --- for a breakthrough is not a heavier step but
  another kind of thing, and ``another kind'' is the very meaning of a
  ratio no counting can reach.\\[0.3em]
  The forces had told us this in numbers; the mind tells us in words; the
  two never spoke, yet say the same.\\[0.5em]
  $\displaystyle \Nstar=\frac{\db}{\dv}\notin\mathbb{Q}
  \quad\Longleftrightarrow\quad
  \text{``of kind, not of magnitude.''}$
}
\end{center}

% Governance Note: This document follows Upgrade Protocol v1.3
% and CRF Style Guide v3.7.


% [BRIDGE-PROSE: integration-added, Style Guide v3.7 §31.6 / UPG-30]
\noindent
With the two kinds shown genuinely distinct, the next question is what they
\emph{build}. ZC91 leaves open (GAP-ZC91-01) how accumulated forward choices
aggregate into a usable quantity. ZC92 answers part of it by constructing the
Barami functional --- a capacity built from the forward-choice subset --- and
identifying a common rate that governs how it grows.
\section{Background: When Algebra Runs Out, Read the Boundary}
\label{sec:background-zc92}
%%=============================================================

%% ── [v3.6 §31.2 / UPG-25] READER'S ON-RAMP ───────
The previous paper showed that the ``breakthrough'' and the ``ordinary
step'' of a cultivating mind can never stand in a clean numerical ratio:
the difference between them is one of kind, and that is exactly what an
irrational ratio means. But it also showed that the ratio is not even a
fixed number --- it depends on how much a mind has already cultivated.
That left an obvious and harder question: as a mind cultivates without
end, what happens to these two kinds of step? CRF's source material gives
two pieces of guidance that turn out to be mathematical, not merely
inspirational. The first, from the Buddha, is that in the limit of deep
cultivation the freedom of choice and the lawfulness of nature stop
disagreeing. The second, from Luang Por Ruesi Ling Dam, is that every act
of choice is gated by three faculties --- noticing, considering, resolving
--- so that the ``cost'' of a forward step is really the inverse of a
cultivated capacity. This paper takes those two statements literally and
finds that, together, they pin down the shape of the answer even though
its exact size remains for a later paper.

%% ── [NEW BLOCK 1] NUMERICAL SUMMARY ────────────────────────
\begin{keybox}[title={Numerical Summary --- ZC92 (T-canonical)}]
{\footnotesize\renewcommand{\arraystretch}{1.3}
\begin{center}
\begin{tabular}{@{}p{5.6cm}p{3.0cm}p{3.6cm}@{}}
\toprule
\textbf{Quantity} & \textbf{Value} & \textbf{Source / Status} \\
\midrule
$\Ninf=\lim_{B\to\infty}\db/\dv$
  & finite, nonzero, irrational
  & \textbf{THM/COR-ZC92 (new)} \\
$\Nstar(B)$ at finite $B$
  & history functional
  & ZC91 (COR-ZC91-02) \\
$\lim_{B\to\infty}(\dv,\db)$
  & $\to 0$ (both)
  & EQ-MB03 (Part~F) \\
common vanishing rate
  & required
  & THM-ZC92-01 \\
value of $\Ninf$
  & undetermined
  & GAP-ZC92-01 (new) \\
\bottomrule
\end{tabular}
\end{center}}
\end{keybox}

%% ── FRAMEWORK CONTEXT ───────────────────────────────────────
\subsection{The Five Constraints and the Two Boundary Records}

The Mind/Barami corpus and ZC91 supply five constraints the functional
must satisfy. EQ-MB01 fixes positivity and ordering ($\db\gg\dv>0$ for all
$B$). EQ-MB03 fixes the asymptotic vanishing ($\dv,\db\to0$ as
$\Barami\to\infty$). The Closure passage of Part~F states that Barami is
\emph{capacity} --- ``a mind with deep barami finds \emph{visesa}
inexpensive'' --- which makes $\dv$ monotone decreasing in $B$. THM-ZC91-01
fixes pointwise irrationality of $\Nstar(B)$ at every $B$. And P-MB04 (the
Buddha's ``inversion at the asymptote'') states that in the limit the
discrete Agentive-Mode equation EQ-MB01 reduces to a gradient flow ---
free will and determinism agree.

Two of these five are doctrinal in origin, and we treat them exactly as
what they are in the corpus: \emph{boundary conditions}. The Buddha's
asymptote is a limit condition on the dynamics. Luang Por Ruesi Ling Dam's
three decision-levels (DEF-MB06: {\thaifont สติ} \emph{sati} / {\thaifont สัมปชัญญะ}
\emph{sampajañña} / {\thaifont อธิษฐาน} \emph{adhiṭṭhāna}) are a gating
condition on the cost. Neither is invoked as authority; each is read for
its mathematical content. This is the standard CRF stance (Map-Not-Reality,
DEF-RN05): the records are data about the structure, not the structure
itself.

\subsection{What Was Already Known Before ZC92}

\begin{keybox}[title={Pre-ZC92 Status}]
\textbf{Known:}
\begin{itemize}[leftmargin=2em, noitemsep]
  \item $\Nstar=\db/\dv$ irrational at every $B$ (THM-ZC91-01); no
    constant value at finite $B$ (COR-ZC91-02).
  \item $\dv,\db\to0$ as $\Barami\to\infty$ (EQ-MB03); $\dv$ monotone
    decreasing (Closure: Barami is capacity).
  \item In the limit, EQ-MB01 $\to$ gradient flow (P-MB04, the Buddha's
    asymptote).
  \item A decision is gated by three levels (DEF-MB06, Luang Por).
\end{itemize}

\textbf{Open (GAP-ZC91-01):}
\begin{itemize}[leftmargin=2em, noitemsep]
  \item The functional form $\dv(B),\db(B)$.
  \item Whether $\Nstar(B)$ has a finite limit, and whether that limit is
    structural.
\end{itemize}

ZC92 pins the joint structure and the asymptotic ratio (THM/COR-ZC92-01)
and the cost mechanism (DEF-ZC92-01); it leaves the explicit form and the
value of $\Ninf$ to GAP-ZC92-01.
\end{keybox}

%% ── [NEW BLOCK 2] SCOPE EXCLUSION ───────────────────────────
\begin{keybox}[title={What ZC92 Does NOT Do}]
\begin{itemize}[leftmargin=2em, noitemsep]
  \item Does not compute $\Ninf$ or the explicit $\dv(B)$ (needs an
    EQ-MB02 Free-Energy calibration; GAP-ZC92-01).
  \item Does not identify $\Ninf$ with the force-layer $36.9638$ or any
    specific number --- that would repeat the cross-layer-equality error
    ZC91 forbade (OBS-ZC92-01).
  \item Does not claim a unique functional \emph{form}; many fit the five
    constraints. Only the common-rate / finite-irrational-limit
    \emph{structure} is pinned.
  \item Does not modify DEF-MB01--07, EQ-MB01--04, or ZC89/90/91.
  \item Does not make a doctrinal claim; it reads the Phase-4 Observer
    record as boundary data and states the mathematics.
\end{itemize}
\end{keybox}

%% ── [NEW BLOCK 3] COMPANION PAPERS ─────────────────────────
\begin{keybox}[title={Companion Papers --- ZC92}]
{\footnotesize\renewcommand{\arraystretch}{1.3}
\begin{center}
\begin{tabular}{@{}p{6.1cm}p{6.4cm}@{}}
\toprule
\textbf{Paper} & \textbf{What it provides to ZC92} \\
\midrule
\texttt{CRF\_ZC91\_IrrationalityOfKind\_v1.tex}
  & THM-ZC91-01, COR-ZC91-02, GAP-ZC91-01 ---
    \textbf{the gap this paper partially closes} \\
\texttt{CRF\_ZC90\_CrossLayerQuantumPair\_v1.tex}
  & DEF-ZC90-01 (Q-pair); CONJ-ZC90-01 (partially rescued here) \\
\texttt{CRF\_Complete\_v17\_16\_Part\_F.tex}
  & EQ-MB02/03, DEF-MB06, P-MB04, Closure; Phase-4 Observer record ---
    \textbf{KEY PILLARS} \\
\texttt{CRF\_Complete\_Index\_v17\_16.tex}
  & Part~F (Mind/Barami) context \\
\bottomrule
\end{tabular}
\end{center}}
\end{keybox}

%% ── DISCOVERY NOTE (promoted) ──
\begin{derivedbox}[title={Discovery Note}]
The session was asked to derive the Barami functional and, on the user's
suggestion, to consult the Phase-4 Observer record. That suggestion was
the turning point. Treated as inspiration the records add nothing to a
proof; treated as boundary conditions they are decisive. The Buddha's
``free will and determinism agree at the limit'' forces a finite
asymptotic ratio (otherwise the jump never merges into the flow); Luang
Por's three levels identify the cost as inverse capacity. The functional
that algebra alone left under-determined became structurally pinned the
moment the doctrinal statements were read as mathematics.
\end{derivedbox}

%% ── PRE-WORK AUDIT ──────────────────────────────────────────
\subsection{Pre-Work Audit (PWA-1 through PWA-10)}

\begin{formalbox}[title={Pre-Work Audit Record --- ZC92}]
Scanned before writing (Upgrade Protocol v1.3, PWA-1 through PWA-10):
\begin{itemize}[leftmargin=2em, noitemsep]
  \item ZC91/THM-ZC91-01, COR-ZC91-02, GAP-ZC91-01 --- the inherited frame.
  \item Part~F/EQ-MB03, P-MB04, Closure passage --- \textbf{KEY PILLARS}
    (vanishing, asymptote, capacity), read verbatim.
  \item Part~F/DEF-MB06 (Three Levels, Luang Por) --- cost-gating mechanism.
  \item ZC90/CONJ-ZC90-01 --- the conjecture partially rescued here.
\end{itemize}

\textbf{PWA-DOWNSTREAM} (ZC81 rule): traced the functional to EQ-MB03 /
P-MB04 / Closure \emph{before} proposing structure; pinned only what the
constraints force, left the value open rather than fitting one.\\
\textbf{Cross-layer-equality guard:} explicitly did \emph{not} set
$\Ninf=36.9638$ (force layer); ZC91 forbade cross-layer numeric identity.
Recorded as OBS-ZC92-01.\\
\textbf{PWA-6} (macro convention): paper macros all
\texttt{\textbackslash providecommand}; no Preamble collision.\\
\textbf{PWA-7/8/9}: not triggered (no R-layer bridge, no new floor).\\
\textbf{PWA-10} (tautology audit): THM-ZC92-01 uses a doctrinal boundary
condition as a genuine constraint on the dynamics, not as a hidden
definition of the result; flagged and discussed in the theorem Remark.\\
\textbf{No duplicate atomic units found.}
\end{formalbox}

\newpage

%%=============================================================
%% ATOMIC UNITS — DEF → THM → COR → FIND → GAP → OBS
%%=============================================================

%%-------------------------------------------------------------
\section{The Cost Mechanism: Capacity Gating\quad(DEF-ZC92-01)}
\label{sec:def-zc92}
%%-------------------------------------------------------------

% [BRIDGE-PROSE: stakes before the box, UPG-30]
Before the limit theorem we need to say, structurally, why the cost of a
forward step should fall as a mind cultivates. The corpus answer is that
Barami is not merely a record but a capacity, and Luang Por's three levels
say what that capacity is made of. We package this as a definition --- not
a closed form, but the named structure the next paper will calibrate.

\needspace{8\baselineskip}
\begin{formalbox}[title={DEF-ZC92-01: Capacity Functional and Cost Law
  ($R_0$, Definition)}]

\textbf{Statement.} Let $\Bcap(B)$ be the \emph{decision-capacity} of a
mind with accumulated Barami $B$, gated multiplicatively by the three
levels of DEF-MB06,
\[
  \Bcap(B)\;\propto\;
  \underbrace{s(B)}_{\text{sati}}\cdot
  \underbrace{p(B)}_{\text{sampajañña}}\cdot
  \underbrace{a(B)}_{\text{adhiṭṭhāna}},
\]
each factor non-decreasing in $B$ (cultivation does not reduce a
faculty). The cost of a forward step is inverse to capacity:
\begin{equation}
  \boxed{\;\dv(B)\;\sim\;\frac{1}{\Bcap(B)},\qquad
  \Bcap(B)\to\infty\ \text{as}\ B\to\infty.\;}
  \label{eq:costlaw-zc92}
\end{equation}

\textbf{Remarks.} (i) This realises the Closure statement ``deep barami
finds \emph{visesa} inexpensive'' and the monotone-decrease constraint
C3. (ii) $\Bcap\to\infty$ gives $\dv\to0$, consistent with EQ-MB03. (iii)
The proportionality is structural; the explicit factors $s,p,a$ require a
dynamical model (GAP-ZC92-01). (iv) When any level vanishes ($s,p$ or
$a\to0$), $\Bcap\to0$ and the mind defaults to Spontaneous Mode ---
exactly DEF-MB06's ``zero sati'' clause.

\textbf{Status:} Definition (structural; not a closed form).
\end{formalbox}

%%-------------------------------------------------------------
\section{The Common-Rate Theorem\quad(THM-ZC92-01)}
\label{sec:thm-zc92}
%%-------------------------------------------------------------

% [BRIDGE-PROSE: the "so what", outsider-readable]
Now the central result. Two facts are in apparent tension. EQ-MB03 says
both kinds of step shrink to nothing as a mind cultivates without end. The
Buddha's asymptote says the discrete, choice-driven trajectory becomes a
smooth, lawful flow in that same limit. The tension is this: if the
breakthrough step shrank faster than the ordinary step, the two would
drift infinitely far apart in relative terms and the ``flow'' would never
form a single coherent grain; if it shrank slower, the breakthrough would
come to dominate and the trajectory would be all jumps. The only way both
records hold at once is for the two steps to shrink together, in lockstep,
so that their ratio settles to a fixed value. That fixed value is the
grain of the limiting flow. The theorem makes this precise; its payoff,
in plain terms, is that the framework's picture of a cultivating mind is
self-consistent only if breakthroughs and ordinary advances keep pace as
they both fade --- and that keeping-pace defines a single asymptotic
number.

\needspace{10\baselineskip}
\begin{formalbox}[title={THM-ZC92-01: Common-Rate Theorem
  ($R_0$, Near-Formal)}]

\textbf{Statement.} Let $\dv(B),\db(B)>0$ satisfy: (C1) $\db\gg\dv$,
(C2) $\dv,\db\to0$ as $B\to\infty$ (EQ-MB03), (C3) $\dv$ monotone
decreasing, (C4) $\Nstar(B)=\db/\dv$ irrational for all $B$
(THM-ZC91-01), and the boundary condition (C5/Buddha, P-MB04) that the
discrete dynamics converge to a gradient flow as $B\to\infty$. Then
$\dv$ and $\db$ vanish at a \emph{common rate}:
\begin{equation}
  \boxed{\;
  \lim_{B\to\infty}\frac{\db(B)}{\dv(B)}=\Ninf,\qquad
  0<\Ninf<\infty.\;}
  \label{eq:thm-zc92}
\end{equation}

\textbf{Derivation.}

\textit{Step 1 (gradient flow requires a coherent grain).} A gradient
flow is a single-scale continuous dynamics: along the conserved direction
$\mathcal{D}_{\rm mind}$ it has, in the limit, one effective step scale,
not two diverging ones. If $\db/\dv\to\infty$ (breakthrough dominates) or
$\db/\dv\to0$ (breakthrough negligible relative to step), the two
generators would belong to different vanishing scales and the limit could
not be a single gradient flow. C5 therefore requires
$\db/\dv\to\Ninf$ with $0<\Ninf<\infty$.

\textit{Step 2 (common rate is equivalent to a finite ratio).} A finite
nonzero limit of $\db/\dv$ is, by definition, the statement that $\dv$ and
$\db$ are asymptotically proportional --- they vanish at the same rate,
$\db(B)=\Ninf\,\dv(B)\,(1+o(1))$. With DEF-ZC92-01, both are
$\sim\Ninf^{?}/\Bcap(B)$ on the same capacity scale.

\textit{Step 3 (consistency with C1--C4).} $\Ninf>1$ is consistent with
C1 ($\db\gg\dv$). The common rate is consistent with C2--C3 (both vanish,
$\dv$ monotone). C4 (pointwise irrationality) is preserved in the limit
under the mild assumption that $\Nstar(B)$ converges (a monotone or
bounded-variation approach suffices); irrationality of the limit is taken
up in COR-ZC92-01.

\hfill (Near-Formal: structure forced by C1--C5; rate function open.)

\textbf{Remark (the PWA-10 point, addressed openly).} The Buddha's
asymptote (C5) is used as a genuine constraint on the limiting dynamics,
not as a covert definition of $\Ninf$. It tells us the limit is a single
gradient flow; from that \emph{dynamical} fact the finite ratio
\emph{follows}. We do not assume the conclusion. Independently, the
force-layer analogue (ZC89) also exhibits two co-vanishing-context steps
with a finite ratio, a structural echo (not an identity --- see Scope).

\physmeaning{As a mind cultivates without end, its ordinary steps and its
breakthroughs both grow small, but they keep a fixed proportion to each
other; that fixed proportion is the texture of the effortless flow the
tradition calls the far shore.}

\textbf{Status:} Near-Formal Theorem (structure pinned; rate function and
value of $\Ninf$ open, GAP-ZC92-01).
\end{formalbox}

%%-------------------------------------------------------------
\section{Corollary: The Asymptotic Ratio Is Irrational\quad(COR-ZC92-01)}
\label{sec:cor-zc92}
%%-------------------------------------------------------------

% [BRIDGE-PROSE]
The limit ratio inherits the character of the finite-$B$ ratios that
approach it, and this is where ZC91 and a rescued piece of ZC90's
conjecture turn out to be compatible rather than contradictory.

\needspace{8\baselineskip}
\begin{formalbox}[title={COR-ZC92-01: $\Ninf$ Irrational; CONJ-ZC90-01
  Partially Rescued ($R_0$, Near-Formal)}]
\textbf{Statement.} $\Ninf=\lim_{B\to\infty}\db/\dv$ is irrational.

\textbf{Derivation.} By THM-ZC91-01, $\Nstar(B)\notin\mathbb{Q}$ for every
finite $B$. By THM-ZC92-01 the sequence $\Nstar(B)$ converges to $\Ninf$.
A convergent sequence of irrationals can in principle have a rational
limit; here it cannot, because a rational $\Ninf=p/q$ would mean the
limiting gradient flow has commensurable jump and step grains --- i.e.\ in
the limit \emph{nibbedha} becomes a finite multiple of \emph{visesa},
reinstating exactly the ``magnitude'' reduction THM-ZC91-01 forbids for
the kind-distinction, which is a property of the choice structure that
\emph{persists} into the limit (the four-fold distinction does not
disappear; P-MB04 says the dynamics become continuous, not that the
kinds merge). Hence $\Ninf\notin\mathbb{Q}$.

\textbf{Relation to ZC90/ZC91.} COR-ZC91-02 retired CONJ-ZC90-01 as a
\emph{finite-$B$ constant}: correct, and unchanged. ZC92 rescues it only
as an \emph{asymptotic} constant $\Ninf$ --- a different object on a
different domain. Both hold: no constant ratio at finite $B$; a finite
irrational ratio in the limit.

\textbf{Status:} Near-Formal Corollary. CONJ-ZC90-01 \textbf{partially
rescued} (limit domain only). Value of $\Ninf$ open (GAP-ZC92-01).
\end{formalbox}

%%-------------------------------------------------------------
\section{The Gift: A Doctrine Becomes a Limit Theorem\quad(FIND-ZC92-01)}
\label{sec:find-zc92}
%%-------------------------------------------------------------

% [BRIDGE-PROSE: outsider-readable "so what"]
The result has a reading that goes beyond the algebra, and it is the one
that makes the whole detour through the Phase-4 records worthwhile. The
oldest tension in any account of mind is between freedom and law: if every
choice is determined, in what sense is it chosen? CRF's source material
asserts that the two agree ``at the limit,'' and that has always sounded
like a graceful evasion. ZC92 gives it a precise and testable meaning.

\needspace{8\baselineskip}
\begin{giftbox}[title={FIND-ZC92-01: Free Will and Determinism as One
  Trajectory at Two Scales (Gift)}]
\textbf{Statement.} The Phase-4 Observer assertion ``free will and
determinism agree at the limit'' (P-MB04) is exactly
\[
  \lim_{B\to\infty}\Nstar(B)=\Ninf\in(0,\infty),\quad\Ninf\notin\mathbb{Q}.
\]
\textbf{Reading.}
\begin{itemize}[leftmargin=2em, noitemsep]
  \item \textbf{Free will} = the finite-$B$ regime: $\Nstar(B)$ is a
    history functional (COR-ZC91-02), genuinely path-dependent; two minds
    in the same present state differ (P-MB02).
  \item \textbf{Determinism} = the limit regime: the trajectory becomes a
    single gradient flow with a fixed grain $\Ninf$ (THM-ZC92-01); the
    history-dependence washes out.
  \item They are the \emph{same} trajectory viewed at two scales of
    cultivation. The ``agreement'' is a limit, not a paradox waved away.
\end{itemize}

\textbf{Why this is a gift.} A doctrinal claim that resisted formalisation
becomes a statement about the convergence of a ratio. It is, in
principle, falsifiable: a model in which the limit ratio diverges, or is
rational, would contradict it.

\physmeaning{Early in cultivation, your choices are truly yours and truly
unpredictable; cultivated without end, the same choosing becomes a smooth
lawful flow --- not because freedom was an illusion, but because freedom
and law were the near and far views of one path.}

\textbf{Status:} Near-Formal Finding (Gift-class). Falsifiable in
principle once a dynamical model fixes $\Ninf$ (GAP-ZC92-01).
\end{giftbox}

%%-------------------------------------------------------------
\section{Open Gap and Provenance Note\quad(GAP-ZC92-01, OBS-ZC92-01)}
\label{sec:gapobs-zc92}
%%-------------------------------------------------------------

% [BRIDGE-PROSE]
What remains open is now a single, sharp, dynamical question, and a clear
note on a temptation we deliberately refused.

\begin{openqbox}[title={GAP-ZC92-01: The Value of $\Ninf$ and the
  Rate Function}]
\textbf{Open item.}
\begin{itemize}[leftmargin=2em, noitemsep]
  \item the explicit capacity factors $s(B),p(B),a(B)$ of DEF-ZC92-01,
    hence the rate function $\dv(B)$;
  \item the numerical value of $\Ninf$, and whether it is a structural
    constant of the mind layer (and if so, what $\Ztf$-data, if any, it
    involves --- to be determined, \emph{not} assumed).
\end{itemize}

\textbf{What closing it would require.} A Free-Energy / $D_{\rm KL}$
calibration of the agentive trajectory via EQ-MB02 (the
$\Sigma$-Lagrange-constraint Lagrangian), solved over a cultivation
history. This is the first mind-layer result requiring a solved dynamical
model rather than the static definitions plus boundary conditions.

\textbf{Target:} ZC93.
\end{openqbox}

\begin{derivedbox}[title={OBS-ZC92-01: The Cross-Layer-Equality Temptation,
  Refused}]
For provenance: it is tempting to set $\Ninf$ equal to the force-layer
ratio $36.9638$ (ZC89), since both layers exhibit a finite ratio between a
continuous and a discontinuous step. ZC91 (Scope; FIND-ZC91-01 Remark)
established that the two layers share \emph{structure}, not \emph{value};
cross-layer numeric identity was explicitly excluded. ZC92 honours that:
$\Ninf$ is left open, not borrowed. Recorded so a later reader sees the
identification was considered and rejected on principle, not overlooked.
\end{derivedbox}

%%=============================================================
%% ATOMIC REGISTRY (MANDATORY)
%%=============================================================
\begin{formalbox}[title={Atomic Registry --- ZC92}]
\begin{center}
\renewcommand{\arraystretch}{1.3}
\small
\begin{tabular}{lp{8.0cm}l}
\toprule
\textbf{ID} & \textbf{Description} & \textbf{Status} \\
\midrule
DEF-ZC92-01
  & Capacity functional $\Bcap(B)$ (DEF-MB06 gating) + cost law
    $\dv\sim1/\Bcap$
  & Definition \\[3pt]
THM-ZC92-01
  & Common-Rate: $\dv,\db$ vanish at a common rate; $\Nstar\to\Ninf$
    finite
  & Near-Formal \\[3pt]
COR-ZC92-01
  & $\Ninf$ irrational; CONJ-ZC90-01 partially rescued (limit domain)
  & Near-Formal \\[3pt]
FIND-ZC92-01
  & Buddha's free-will/determinism agreement $=$ $\lim\Nstar=\Ninf$ finite
  & Near-Formal (Gift) \\[3pt]
GAP-ZC92-01
  & Value of $\Ninf$; rate function $\dv(B)$; target ZC93
  & Open \\[3pt]
OBS-ZC92-01
  & Cross-layer-equality ($\Ninf=36.9638$) considered and refused
  & Observation \\[3pt]
\midrule
\multicolumn{3}{l}{\emph{Gap-status changes from prior papers:}} \\
GAP-ZC91-01 & Barami-history functional; partially closed via THM-ZC92-01
  & Partial $\checkmark$ \\
CONJ-ZC90-01 & ratio constant; partially rescued on limit domain
  (COR-ZC92-01) & Partial $\checkmark$ \\
\bottomrule
\end{tabular}
\end{center}
\end{formalbox}

%%=============================================================
\section{R-Layer Research Note}
\label{sec:rlayer-zc92}
%%=============================================================
\begin{layerbox}[title={R-Layer Assignments for ZC92}]

\textbf{Layer assignments:}

{\small\renewcommand{\arraystretch}{1.25}
\begin{center}
\begin{tabular}{@{}llll@{}}
\toprule
Atomic unit & R-layer & Cost $T$ & Source \\
\midrule
DEF-ZC92-01 & $R_0$ & $0$ & structural (DEF-MB06, Closure) \\
THM-ZC92-01 & $R_0$ & $0$ & limit argument (EQ-MB03 + P-MB04) \\
COR-ZC92-01 & $R_0$ & $0$ & follows THM-ZC91-01 + THM-ZC92-01 \\
FIND-ZC92-01 & $R_0$ & $0$ & interpretation of P-MB04 \\
GAP-ZC92-01 & $R_0{\to}$ dyn. & --- & EQ-MB02 calibration (not a bridge) \\
\bottomrule
\end{tabular}
\end{center}}

\textbf{Two-layer floor (PWA-8):}
\begin{itemize}[leftmargin=2em, noitemsep]
  \item Spectral layer ($R_0$): $\epscan=0.007550$ --- not used.
  \item Bridge layer: $\epsuniv=0.007599$ --- not used; ZC92 introduces no
    floor. The force-layer ratio is cited only as a structural echo
    (OBS-ZC92-01), not numerically imported.
\end{itemize}

\textbf{Key R-layer remark.} All results at $R_0$, structural / asymptotic;
no Born-chain bridge cost. The open item GAP-ZC92-01 is dynamical, not an
R-layer crossing.
\end{layerbox}

%%=============================================================
\section{Genealogy Map}
\label{sec:genealogy-zc92}
%%=============================================================
\begin{genealogybox}[title={How ZC92's Key Objects Trace Back}]

\textbf{Branch A --- The functional, from definition to structure:}
\begin{itemize}[leftmargin=2em, noitemsep]
  \item \textbf{Part~F / EQ-MB01}: the step magnitudes $\dv,\db$.
  \item $\to$ \textbf{Part~F / EQ-MB03 + Closure}: both vanish; Barami is
    capacity.
  \item $\to$ \textbf{ZC91 / GAP-ZC91-01}: the functional re-opened.
  \item $\to$ \textbf{ZC92 / DEF-ZC92-01, THM-ZC92-01}: cost law +
    common-rate structure.
\end{itemize}

\textbf{Branch B --- The boundary records, from doctrine to mathematics:}
\begin{itemize}[leftmargin=2em, noitemsep]
  \item \textbf{Part~F / P-MB04 (Buddha)}: free will and determinism
    agree at the asymptote.
  \item \textbf{Part~F / DEF-MB06 (Luang Por)}: three decision-levels
    gate the choice.
  \item $\to$ \textbf{ZC92 / THM-ZC92-01}: the asymptote becomes a
    finite-ratio condition; the three levels become the capacity gate.
  \item $\to$ \textbf{ZC92 / FIND-ZC92-01}: the doctrine becomes a limit
    theorem.
\end{itemize}

\textbf{Branch C --- The ratio across three papers:}
\begin{itemize}[leftmargin=2em, noitemsep]
  \item \textbf{ZC90 / CONJ-ZC90-01}: ratio conjectured a structural
    constant.
  \item \textbf{ZC91 / COR-ZC91-02}: no constant at finite $B$.
  \item $\to$ \textbf{ZC92 / COR-ZC92-01}: a constant in the limit,
    irrational --- both prior results reconciled by domain.
\end{itemize}

\textbf{Convergence at ZC92:}
Branch~A could only pin structure; Branch~B's boundary records pinned the
asymptotic ratio; Branch~C shows the three-paper arc on the ratio is
consistent once finite-$B$ and limit domains are separated.

ZC92's contribution: it partially closes the Barami functional by reading
the Phase-4 Observer record as boundary data, and turns ``free will and
determinism agree at the limit'' into a precise, in-principle-falsifiable
limit theorem --- something no static-definition paper could reach.
\end{genealogybox}

%%=============================================================
\section{Closing Sentence}
%%=============================================================
%% UPG-22 — TONE B (Contemplative): partial closure; structure and
%% asymptotic ratio pinned, value open. Declared in header. Precedent: ZC72,
%% ZC75 (algebraic/structural settled; the dynamical / evolution part open).

\bigskip
\begin{center}
\textit{%
  The ordinary step and the breakthrough both grow small as the path is
  walked, yet they fade together, keeping one proportion --- and that
  proportion is the grain of the flow that the choosing finally becomes.\\[0.3em]
  Freedom near, law far; one path. The number that names the grain
  belongs to ZC93.
}
\end{center}

% Governance Note: This document follows Upgrade Protocol v1.3
% and CRF Style Guide v3.7.


% [BRIDGE-PROSE: integration-added, Style Guide v3.7 §31.6 / UPG-30]
\noindent
The Barami functional says capacity accumulates, but not what \emph{controls}
whether a given choice can take the smooth step or the sudden one. ZC92 leaves
this open (GAP-ZC92-01). ZC93 supplies the gate: a ratio, built from the three
faculties, that decides which kind of step a choice is permitted to make ---
partially closing the gap and setting up what the next paper must finish.
\section{Background: A Form Before a Number}
\label{sec:background-zc93}
%%=============================================================

%% ── READER'S ON-RAMP (v3.6 §31.2 / UPG-25) ────────────────
\begin{keybox}[title={If you are just arriving --- ZC93 in plain language}]
\sloppy
CRF treats a mind as something that, at each conscious moment, chooses one of
four moves; two of them carry it \emph{forward} --- an ordinary advance
(\emph{visesa}) and a decisive breakthrough (\emph{nibbedha}). Each has a
``cost''. The question this paper inherits is: as a mind practises without
end, do these two costs settle into a fixed proportion? The earlier papers
said the proportion is irrational (the two moves differ in \emph{kind}, not
\emph{size}) and that it has no fixed value while practice is still finite ---
yet that it does settle to a fixed value in the limit. What was missing was
that value. This paper shows \emph{why the value has a clean form even though
we cannot yet name it}: the deep ``capacity'' a practised mind builds up makes
\emph{both} steps cheaper in the same way, so it cancels out of their ratio.
What is left is a ratio of three faculties --- noticing, considering,
resolving --- and that ratio, not the capacity, is the number. We pin the
form; we leave the number, honestly, for a measurement we do not yet have.
\end{keybox}

%% ── NUMERICAL SUMMARY ──────────────────────────────────────
\begin{keybox}[title={Numerical Summary --- ZC93 (T-canonical)}]
{\footnotesize\renewcommand{\arraystretch}{1.3}
\begin{center}
\begin{tabular}{@{}p{5.8cm}p{3.0cm}p{3.4cm}@{}}
\toprule
\textbf{Quantity} & \textbf{Value} & \textbf{Source / Status} \\
\midrule
$\Ninf=\lim_{B\to\infty}\db/\dv$
  & $=\ab/\av=\Gb/\Gv$
  & \textbf{THM/COR-ZC93 (new)} \\
capacity $\Bcap(B)$ in the ratio
  & cancels exactly
  & THM-ZC93-01 \\
$\Ninf$ : finite, nonzero
  & yes
  & THM-ZC92-01 \\
$\Ninf$ : irrational
  & yes
  & THM-ZC91-01 (inherited) \\
value of $\Gb/\Gv$
  & undetermined
  & GAP-ZC93-01 (new) \\
force-layer analogue $36.9638$
  & \emph{not} equal
  & ZC89 (cited, never identified) \\
\bottomrule
\end{tabular}
\end{center}}
\end{keybox}

\subsection{Where ZC92 Stopped, and Why}

ZC92 reached an unusual situation. It had \emph{proved} that an object exists
--- the asymptotic ratio $\Ninf$ --- and had pinned three of its properties
(finite, nonzero, irrational), yet could not say what it equals. The
obstruction was named plainly: computing $\Ninf$ appeared to require an
EQ-MB02 Free-Energy calibration, a dynamical input fixing the actual
magnitudes of the forward steps as functions of cultivation. The corpus
supplies the \emph{form} of that constraint (EQ-MB02 is a constrained
Lagrangian with an active suffering term) but not a \emph{calibration} of it.

That is a real wall against the \emph{number}. It is not, this paper argues, a
wall against the \emph{form}. The reverse-inference move --- start from what
the limit must be (finite, nonzero, irrational) and ask what that forces ---
shows that the very feature making $\Ninf$ hard to compute (both steps depend
on an uncalibrated capacity) is also what makes its form clean: a quantity
that divides both numerator and denominator cannot survive in their ratio.

\subsection{The Three Constraints This Paper Uses}

We use exactly three inherited results, none modified:
\begin{enumerate}[leftmargin=2em, noitemsep]
  \item \textbf{THM-ZC92-01} (common rate): $\dv(B)\to0$ and $\db(B)\to0$ as
    $\Barami\to\infty$, at a common rate; the limit $\Ninf$ exists, finite and
    nonzero.
  \item \textbf{THM-ZC91-01} (irrationality of kind): $\Nstar(B)$ is
    irrational at every $B$; $\db/\dv$ is never a rational number.
  \item \textbf{DEF-MB06} (three faculties): a forward choice is gated by
    \emph{sati} (noticing), \emph{sampaja\~n\~na} (considering), and
    \emph{adhi\d{t}\d{t}h\=ana} (resolving); absence of any one removes the
    choice from the available set.
\end{enumerate}
The Closure passage of Part~F adds the reading that makes these click
together: \emph{Barami is also capacity} --- ``a mind with deep barami finds
\emph{visesa} inexpensive''. Capacity is exactly the object that will cancel.

%%=============================================================
\section{The Capacity-Cancellation}
\label{sec:cancellation-zc93}
%%=============================================================

THM-ZC92-01 says both forward costs vanish at a \emph{common} rate. The phrase
``common rate'' is the whole content of this section: it means the two costs
are, to leading order in the limit, the \emph{same} decaying function of
cultivation, scaled by two different constants. Write that decaying function as
$1/\Bcap(B)$, where $\Bcap(B)$ is the cultivated capacity (DEF-ZC92-01:
$\dv\sim1/\Bcap$). Then the common-rate statement \emph{is} the statement that
the breakthrough cost shares the same $\Bcap$-scaling, only with its own
constant. The proportion between the two is therefore carried entirely by the
two constants, never by $\Bcap$. The capacity decides how expensive practice
is; it does not decide the shape of the breakthrough relative to the ordinary
step.

\begin{formalbox}[title={THM-ZC93-01: Capacity-Cancellation
  (Exact; closes the structural half of GAP-ZC92-01)}]
\textbf{Hypotheses.} (i) THM-ZC92-01: $\dv,\db\to0$ as $B\to\infty$ at a
common rate, with $\Ninf=\lim_{B\to\infty}\db/\dv$ finite and nonzero.
(ii) DEF-ZC92-01: there is a cultivated capacity $\Bcap(B)$, monotone
increasing in $B$, with $\dv(B)\sim\av/\Bcap(B)$ as $B\to\infty$ for a
constant $\av>0$.

\textbf{Claim.} The common-rate hypothesis is equivalent to
$\db(B)\sim\ab/\Bcap(B)$ for a constant $\ab>0$ \emph{with the same}
$\Bcap(B)$, and then
\[
  \Ninf=\lim_{B\to\infty}\frac{\db(B)}{\dv(B)}
       =\lim_{B\to\infty}\frac{\ab/\Bcap(B)}{\av/\Bcap(B)}
       =\frac{\ab}{\av},
\]
in which $\Bcap(B)$ cancels exactly. Hence $\Ninf$ is a structural constant
(a ratio of two leading coefficients) although neither $\dv$ nor $\db$ is
constant in $B$.

\textbf{Proof.} If $\db$ vanished at a \emph{different} rate from $\dv$, the
ratio $\db/\dv$ would tend to $0$ or $\infty$, contradicting the finite,
nonzero limit of hypothesis~(i). So $\db$ and $\dv$ have the same leading
decay; writing $\dv\sim\av/\Bcap$ from~(ii), the common leading factor is
$1/\Bcap$, and $\db\sim\ab/\Bcap$ with $\ab=\Ninf\,\av$. The displayed limit
follows by direct substitution; $\Bcap(B)$ appears identically in numerator
and denominator and cancels. $\qquad\blacksquare$
\end{formalbox}

The unboxed point worth keeping is the one an outsider can carry away without
the algebra: \emph{the thing that is hard to measure --- how much capacity a
mind has built --- is exactly the thing that does not affect the answer.} The
proportion of breakthrough to ordinary step is a property of the steps
themselves, read in the limit, not of how cheaply a given mind can pay for
them. This is why ZC91's ``no fixed value at finite $B$'' and ZC92's ``a fixed
value in the limit'' are not in tension: at finite cultivation the capacity is
still entangled with history (P-MB02), and the clean cancellation has not yet
been reached; only in the limit do the two costs become the same decaying
function, and only then does $\Bcap$ drop out.

\physmeaning{What a mind builds through practice (capacity) changes the price
of every forward step but not the ratio between a breakthrough and an ordinary
advance. That ratio is fixed by the kind of steps, not by the practitioner's
accumulated ease.}

%%=============================================================
\section{The Surviving Coefficients Are Gates}
\label{sec:gates-zc93}
%%=============================================================

THM-ZC93-01 leaves $\Ninf=\ab/\av$: a ratio of two leading coefficients with
no capacity in sight. The next question is what those coefficients \emph{are}.
Here the second piece of the Phase-4 record does the work that algebra alone
cannot. Luang Por Ruesi Ling Dam's decomposition of every decision into three
faculties (DEF-MB06) is not a piece of devotional framing; it is the cost
mechanism. A forward choice is \emph{available} only when all three faculties
are present, and the ease with which it is taken is set by how fully each is
present. That product of availabilities is what we name a \emph{gate}.

\begin{formalbox}[title={DEF-ZC93-01: The Faculty-Gate}]
For a forward choice $c\in\Cfourplus=\{\text{visesa},\text{nibbedha}\}$, let
$\Fsati,\Fsamp,\Fadh\in(0,1]$ be the availabilities at moment $n$ of the three
DEF-MB06 faculties --- \emph{sati} (noticing the moment as a decision-point),
\emph{sampaja\~n\~na} (considering the available choices), and
\emph{adhi\d{t}\d{t}h\=ana} (placing the resolution). The \emph{faculty-gate}
of $c$ is the product
\[
  \Gate{c}=\Fsati_{c}\,\Fsamp_{c}\,\Fadh_{c}\,,
\]
and the cultivated capacity entering DEF-ZC92-01 is the gate-weighted record
$\Bcap(B)=\Bcap\big(\{\Gate{c_n}\}_{n\le N}\big)$ accumulated over the
forward choices of DEF-MB04. A faculty at zero closes the gate
($\Gate{c}=0$): the choice leaves the available set, recovering the DEF-MB06
statement that absence of any one faculty removes the choice. ``Saturation''
is the limit $\Fsati,\Fsamp,\Fadh\to1$, in which the gate reaches its
structural maximum $\Gate{c}\to\Gate{c}^{\,\mathrm{sat}}$.
\end{formalbox}

The leading coefficients $\av,\ab$ of THM-ZC93-01 are read off at saturation:
in the deep-cultivation limit the three faculties are fully available (this is
what ``deep barami'' \emph{means} operationally --- the mind reliably notices,
considers, and resolves), so the residual difference between the ordinary and
the breakthrough step is the difference between their \emph{saturated gates}.
The ordinary step passes the ordinary forward gate $\Gv=\Gate{v}^{\,
\mathrm{sat}}$; the breakthrough passes the penetrative gate $\Gb=\Gate{b}^{\,
\mathrm{sat}}$. These are not equal, because nibbedha is a discontinuous
penetration and visesa is a continuous advance (DEF-MB03: a difference of kind,
not of magnitude); and being of kind, their ratio cannot be rational
(THM-ZC91-01).

\begin{formalbox}[title={COR-ZC93-01: Gate-Ratio Form of the Asymptotic
  Constant}]
Combining THM-ZC93-01 with DEF-ZC93-01 at saturation,
\[
  \Ninf=\frac{\ab}{\av}=\frac{\Gb}{\Gv}
       =\frac{\Gate{b}^{\,\mathrm{sat}}}{\Gate{v}^{\,\mathrm{sat}}}\,,
\]
the ratio of the saturated nibbedha and visesa faculty-gates. By
THM-ZC91-01 the ratio is irrational; by THM-ZC92-01 it is finite and nonzero.
Thus $\Ninf$ is a \emph{gate ratio}: a pure structural proportion between two
kinds of forward choice, independent of the capacity $\Bcap$ that sets their
absolute cost.
\end{formalbox}

\physmeaning{In a fully cultivated mind, ordinary advance and decisive
breakthrough are both freely available; what still distinguishes them is the
\emph{kind} of gate each passes through. The asymptotic ratio measures that
difference of kind --- and because it is a difference of kind, it never
reduces to a whole-number multiple.}

%%=============================================================
\section{Phase-4 Observer: Two Records, One Boundary Condition}
\label{sec:phase4-zc93}
%%=============================================================

This framework's mind layer was begun from outside academic physics, and its
author has been candid that the language available to describe it was thinner
than the structure it pointed at. The Phase-4 Observer track exists for
precisely that reason: the source records are read not as inspiration around
the mathematics but as \emph{boundary data} inside it. ZC92 established the
precedent --- the Buddha's asymptote was the boundary condition that turned an
under-determined functional into a pinned structure. This paper makes the
pairing explicit, because the two records turn out to answer two different
halves of the same question.

\begin{genealogybox}[title={How the two records translate}]
\textbf{Record~1 --- the Buddha (P-MB04, CLM-MB03).}
``As barami $\to\infty$ the Agentive-Mode equation reduces to a gradient flow;
free will and determinism agree at the limit.'' Read as mathematics: the limit
$B\to\infty$ is \emph{well-defined} and the discrete forward dynamics converge
to a single continuous flow. This is exactly the guarantee that
$\lim_{B\to\infty}\db/\dv$ \emph{exists} as a finite, nonzero number --- the
existence half. Without it, $\Ninf$ would be a symbol for a limit that might
not converge.

\smallskip
\textbf{Record~2 --- Luang Por Ruesi Ling Dam (DEF-MB06).}
Every decision is gated by three faculties --- noticing, considering,
resolving. Read as mathematics: the cost of a forward step factorises through
a product of three availabilities (DEF-ZC93-01), and at saturation the
residual proportion between two kinds of forward step is a ratio of their
gates. This is exactly \emph{what the limit is a ratio of} --- the content
half.

\smallskip
\textbf{Together.} The first record says the limit \emph{exists}; the second
says \emph{what it is a ratio of}. Neither alone pins the structure; together
they are a single boundary condition: \emph{$\Ninf$ is the well-defined
saturation gate-ratio $\Gb/\Gv$.} A pair of doctrinal statements becomes one
boundary condition for a limit theorem.
\end{genealogybox}

\begin{giftbox}[title={FIND-ZC93-01: The Two Records Are One Boundary Condition}]
The statement ``free will and determinism agree at the limit'' (P-MB04) and
the three faculties of decision (DEF-MB06) are not two separate inputs to this
calculation; they are the existence half and the content half of a single
boundary condition on $\Ninf$. The Buddha's asymptote guarantees the limit
converges; Luang Por's faculties identify the converged value as a gate ratio.
Read this way, a doctrinal pair becomes the boundary data of a limit theorem
--- extending FIND-ZC92-01 (``the Buddha's agreement \emph{is} the limit
statement'') from existence to form. The mathematics did not need new
spiritual language; it needed the records read as constraints, and they were
sufficient.
\end{giftbox}

There is a quiet structural reward here for the framework's founding wager ---
that a physics able to touch the mind layer would, in return, give that
otherwise-unmeasurable layer some independent weight. ZC93 is a small instance
of the wager paying off in the expected direction: a claim that lives in the
contemplative tradition (agreement at the limit; choice gated by faculties)
has been made to carry the existence and the form of a quantity in a derivation
chain, with no doctrinal special-pleading. The records earned their place by
functioning as boundary conditions a calculation could not do without --- not
by assertion, but by use.

%%=============================================================
\section{Scope and the Forbidden Identification}
\label{sec:scope-zc93}
%%=============================================================

The force layer carries an exactly analogous asymptotic constant: the
irrational step ratio $36.9638\ldots$ (FIND-ZC89-01). It is tempting, and it
would be wrong, to set $\Ninf$ equal to it. ZC91's FIND established that the
two layers are forced irrational by \emph{independent} logic --- the force
layer by ladder arithmetic, the mind layer by the kind/magnitude distinction
--- and ZC92's Scope forbade identifying their values. ZC93 honours that
boundary: the force-layer number is cited here only as the structurally
parallel object, never as the value of $\Gb/\Gv$. The parallel is that both
are gate-like irrational ratios of two non-commensurate generators; the
non-identity is that nothing in the corpus relates their magnitudes, and a
great deal (FIND-ZC91-01) says their necessities are unrelated.

\begin{openqbox}[title={OBS-ZC93-01: The cross-layer-equality temptation,
  pre-empted}]
Probe considered and declined \emph{before} drafting: ``$\Ninf=36.9638$, the
force-layer ratio.'' Declined for lack of any structural relation between the
mind-layer gate magnitudes and the force-layer $\mathbb{Z}_{15}$ ladder data;
asserting equality would repeat exactly the error ZC91--ZC92 ruled out. No
scar is registered because the probe was never run as a claim --- recorded as
an observation for provenance, consistent with the ZC89-E4 discipline against
fabricated or wishful numbers.
\end{openqbox}

%%=============================================================
\section{Open Gap}
\label{sec:gap-zc93}
%%=============================================================

The honest residue is small and sharp. The \emph{form} of $\Ninf$ is settled
($=\Gb/\Gv$, a saturated gate ratio, irrational, finite, nonzero). The bare
\emph{number} is not, and cannot be without an input the corpus does not yet
contain: a calibration of the saturated gate magnitudes themselves.

\begin{openqbox}[title={GAP-ZC93-01: The Value of the Gate Ratio
  (Priority~3; future EQ-MB02 calibration)}]
\textbf{Statement.} What is the numeric value of $\Gb/\Gv$, the ratio of the
saturated nibbedha and visesa faculty-gates? COR-ZC93-01 fixes that it is
irrational, finite, and nonzero, but not its magnitude.

\textbf{Why it is distinct.} THM-ZC93-01 and COR-ZC93-01 are exact and closed:
the capacity cancels and the surviving object is a gate ratio. This gap is the
\emph{calibration} question only --- deliberately separated from the structural
result so the latter is not held hostage to the former (the same discipline
ZC88 used in separating its structural bridge from the metric step-ratio).

\textbf{Approach candidates.}
\begin{enumerate}[leftmargin=2em, noitemsep]
  \item An EQ-MB02 Free-Energy calibration: instantiate the constrained
    Lagrangian $L_{\rm mind}=T_{\rm kin}-V_{\rm pot}+\lambda(\Sigma-
    \Sigma_{\rm crit})_+$ with a measured or modelled $\lambda$ and extract the
    saturated gate magnitudes.
  \item An empirical decision-cost programme operationalising DEF-MB06 (rates
    of noticing / considering / resolving) to estimate $\Gb,\Gv$ directly.
\end{enumerate}
\textbf{Status:} Open. Target: a future EQ-MB02 calibration, \emph{not} the
next paper. The number is deferred, not assumed.
\end{openqbox}

%%=============================================================
\section{Master Status Table}
%%=============================================================

{\small\renewcommand{\arraystretch}{1.3}
\begin{longtable}{@{}p{4.8cm}p{2.4cm}p{5.0cm}@{}}
\toprule
\textbf{Atomic unit} & \textbf{Status} & \textbf{Note} \\
\midrule
\endhead
THM-ZC93-01 ($\Bcap$-cancellation) & Exact
  & capacity divides out; $\Ninf=\ab/\av$ \\[4pt]
COR-ZC93-01 (gate-ratio form) & Exact
  & $\Ninf=\Gb/\Gv$; finite, nonzero, irrational \\[4pt]
DEF-ZC93-01 (faculty-gate) & Definition
  & $\Gate{c}=\Fsati\Fsamp\Fadh$; capacity is gate-weighted \\[4pt]
FIND-ZC93-01 (records = one BC) & Exact (Gift)
  & Buddha = existence, Luang Por = content \\[4pt]
GAP-ZC92-01 (value of $\Ninf$) & Structural half CLOSED
  & form pinned by THM/COR-ZC93; number re-opened \\[4pt]
GAP-ZC93-01 (gate-ratio value) & Open
  & needs EQ-MB02 calibration; future, not ZC94 \\[4pt]
OBS-ZC93-01 (cross-layer eq.) & Observation
  & equality with 36.9638 declined pre-emptively \\
\midrule
\multicolumn{3}{@{}p{12.5cm}}{\textit{Gap-status changes from prior papers:}
  GAP-ZC92-01 structural half closed (form of $\Ninf$ pinned: gate ratio,
  capacity cancels). No value fabricated. ZC89/90/91/92 and all DEF-MB /
  EQ-MB units used verbatim as pillars; none modified.} \\
\bottomrule
\end{longtable}
}

%%=============================================================
%% ATOMIC REGISTRY (MANDATORY)
%%=============================================================

\begin{formalbox}[title={Atomic Registry --- ZC93}]
\begin{center}
\renewcommand{\arraystretch}{1.3}
\small
\begin{tabular}{lp{7.8cm}l}
\toprule
\textbf{ID} & \textbf{Description} & \textbf{Status} \\
\midrule
THM-ZC93-01
  & Under common-rate vanishing, capacity $\Bcap$ cancels in $\db/\dv$;
    $\Ninf=\ab/\av$
  & Exact \\[3pt]
COR-ZC93-01
  & $\Ninf=\Gb/\Gv$, the saturated faculty-gate ratio; finite, nonzero,
    irrational
  & Exact \\[3pt]
DEF-ZC93-01
  & Faculty-gate $\Gate{c}=\Fsati\Fsamp\Fadh$; $\Bcap$ is gate-weighted
    record
  & Definition \\[3pt]
FIND-ZC93-01
  & P-MB04 (existence) + DEF-MB06 (content) are one boundary condition on
    $\Ninf$
  & Exact (Gift) \\[3pt]
GAP-ZC93-01
  & Numeric value of $\Gb/\Gv$; needs EQ-MB02 calibration
  & Open \\[3pt]
OBS-ZC93-01
  & Cross-layer equality $\Ninf=36.9638$ declined (no structural relation)
  & Observation \\
\midrule
\multicolumn{3}{@{}p{11.5cm}}{\textit{Inherited verbatim (not new):}
  THM-ZC92-01, DEF-ZC92-01, FIND-ZC92-01 (ZC92); THM-ZC91-01, COR-ZC91-02
  (ZC91); DEF-MB03/04/06, EQ-MB01/02/03/04, P-MB04, CLM-MB03 (Part~F);
  FIND-ZC89-01 (force-layer analogue, cited only). No value modified.} \\
\bottomrule
\end{tabular}
\end{center}
\end{formalbox}

%%=============================================================
\section{R-Layer Research Note}
\label{sec:rlayer-zc93}
%%=============================================================

\begin{layerbox}[title={R-Layer Assignments for ZC93}]

\textbf{Layer assignments:}

{\small\renewcommand{\arraystretch}{1.25}
\begin{center}
\begin{tabular}{@{}llll@{}}
\toprule
Atomic unit & R-layer & Cost $T$ & Source \\
\midrule
THM-ZC93-01  & $R_0$ & $0$ & limit algebra from THM-ZC92-01 + DEF-ZC92-01 \\
COR-ZC93-01  & $R_0$ & $0$ & DEF-ZC93-01 at saturation + THM-ZC91-01 \\
DEF-ZC93-01  & $R_0$ & $0$ & structural definition (DEF-MB06 factorisation) \\
FIND-ZC93-01 & $R_0$ & $0$ & reading of P-MB04 + DEF-MB06 as one BC \\
\bottomrule
\end{tabular}
\end{center}}

All ZC93 results are at the $R_0$ layer: no Born-chain bridge running enters,
and no spectral-layer floor value is touched. The result is a limit identity
plus a definitional factorisation of an already-established cost mechanism
(DEF-ZC92-01); it adds no new dynamical content, only the cancellation that
the common-rate hypothesis already implies.
\end{layerbox}

%%=============================================================
\section{Genealogy Map}
\label{sec:genealogy-zc93}
%%=============================================================

\begin{genealogybox}[title={How ZC93's Key Objects Trace Back}]

\textbf{Branch --- the mind-layer ratio:}
\begin{itemize}[leftmargin=2em, noitemsep]
  \item \textbf{DEF-MB03/04 (Part~F)}: the forward pair $\Cfourplus=\{$visesa,
    nibbedha$\}$; visesa--nibbedha is a difference of kind.
  \item $\to$ \textbf{ZC90/DEF-ZC90-01}: the pair is a Q-pair (one conserved
    direction, two non-commensurate generators).
  \item $\to$ \textbf{ZC91/THM-ZC91-01}: the ratio $\db/\dv$ is irrational at
    every $B$; no finite-$B$ constant (COR-ZC91-02).
  \item $\to$ \textbf{ZC92/THM-ZC92-01}: common-rate vanishing; the limit
    $\Ninf$ exists, finite, nonzero, irrational. Value left open
    (GAP-ZC92-01).
  \item $\to$ \textbf{ZC93/THM-ZC93-01}: capacity $\Bcap$ cancels in the
    limit; $\Ninf=\ab/\av$.
  \item $\to$ \textbf{ZC93/COR-ZC93-01}: $\ab/\av=\Gb/\Gv$, the saturated
    faculty-gate ratio (DEF-MB06).
\end{itemize}

\textbf{Boundary data (Phase-4 Observer record, Part~F):}
\begin{itemize}[leftmargin=2em, noitemsep]
  \item \textbf{P-MB04 / CLM-MB03} (the Buddha): the limit exists; free will
    and determinism agree there.
  \item \textbf{DEF-MB06} (Luang Por Ruesi Ling Dam): the three faculties gate
    each decision; the surviving coefficients are gates.
\end{itemize}

\textbf{Convergence at ZC93:}
The mind-layer branch supplies an asymptotic constant whose existence ZC92
proved but whose value seemed to need a calibration. The capacity-cancellation
removes the calibration from the \emph{form}: what is left is a ratio of two
saturated faculty-gates. The Phase-4 record supplies both the existence (Buddha)
and the content (Luang Por) of that ratio. ZC93's contribution: it pins the
form of $\Ninf$ as a gate ratio, draws the exact boundary against identifying
it with the force-layer $36.9638$, and defers only the bare number, honestly,
to a future calibration.
\end{genealogybox}

%%=============================================================
\section{Closing Sentence}
%%=============================================================
%% TONE B — CONTEMPLATIVE (structural half closed; value left open).

\bigskip
\begin{center}
\textit{%
  The deeper a mind is cultivated, the cheaper every forward step becomes ---
  and in that cheapening, the cost itself falls away, leaving only the bare
  proportion between an ordinary advance and a breakthrough.\\[0.3em]
  That proportion has a shape now; its size waits for a measure we do not yet
  hold.
}

\medskip
{\thaifont
เมื่อบารมีลึกพอ ทุกก้าวไปข้างหน้าก็เบาลงพร้อมกัน ---
สิ่งที่เคยเป็นราคาของแต่ละก้าวจึงหักล้างกันไป เหลือไว้เพียงอัตราส่วน
ระหว่างการก้าวอย่างสามัญกับการก้าวที่ทะลุ ---
อัตราส่วนที่ไม่มีวันลงตัว ดั่งที่ทรงตรัสว่า
ณ ที่สุดนั้น เจตจำนงและกฎเกณฑ์มิได้ขัดแย้งกันเลย.}
\end{center}

% Governance Note: This document follows Upgrade Protocol v1.3
% and CRF Style Guide v3.7 (UPG-30 prose floor, UPG-31 header contrast).
% Gauge: T-canonical. Tone B. Phase-4 Observer track active.


%%=============================================================
\part{Cross-Part Connection Map v17.16 (Part L)}
\label{part:LIX}
%%=============================================================

\begin{conmapbox}[title={Cross-Part Connection Map v17.16 --- Reading
  Conventions}]
\small
This map records, for each paper in Part L (ZC88--93):
(1) \textbf{Backward Inheritance} --- which objects from earlier Parts are
used as input; (2) \textbf{Forward Corrections} --- which prior gaps are
closed or refined; (3) \textbf{Internal Gap Chain} --- how the six papers
hand gaps to one another. Source files are never modified (No-Loss).

\textbf{Notation:}
\textcolor{crfgreen}{$\checkmark$ CLOSED} = gap resolved;
\textcolor{crfgreen}{$\approx$ PARTIAL} = gap partially closed;
$\star$ RETIRED = conjecture withdrawn.
\textbf{PM-V17\_16-EPSFAMILY:} local overrides
$\backslash$\texttt{epsEM}$\to\varepsilon_{\rm EM}$,
$\backslash$\texttt{epscan}$\to\varepsilon_{\rm can}$,
$\backslash$\texttt{epsSC}$\to\varepsilon_{\rm SC}$,
$\backslash$\texttt{epsS}$\to\varepsilon_{\rm strong}$
(bare sector forms, Style Guide v3.7); continues PM-V17\_15-EPSFAMILY of
Part K. Chain and sources unmodified.
\end{conmapbox}

\begin{inheritbox}[title={Backward Inheritance --- Part L}]
\small
\textbf{Cluster 1 (Force-Layer, ZC88--89)} inherits from \textbf{Part K}
(ZC83--87): the floor lattice and its multiplicative (geometric-mean)
ordering, the $\Gamma$-ladder and its additive ordering, and the
$\varepsilon$-family values (verified-consistent; see Part K's
do-not-reconcile directive, which remains in force here).

\medskip
\textbf{Cluster 2 (Mind-Layer, ZC90--93)} inherits from \textbf{Part F}
(ZC41--44): the four-way decision structure $\mathcal{C}_4$ (DEF-MB03) and
its forward subset $\mathcal{C}_4^{+}$ (DEF-MB04); the Barami object
(DEF-MB04/05) and the three faculties \textit{sati / sampaja\~n\~na /
adhi\d{t}\d{t}h\=ana} (DEF-MB06); the decision-driven trajectory and its
Lagrange form (EQ-MB01/02); the visesa/nibbedha step pair; and the
asymptotic inversion principle (P-MB04). ZC90--93 add structure on top of
these pillars and change no Part F definition.
\end{inheritbox}

\begin{correctionbox}[title={Forward Corrections and Internal Gap Chain --- Part L}]
\small
The six papers form a single continuous gap chain:
\begin{itemize}[noitemsep, leftmargin=1.6em]
  \item \textbf{ZC88} opens \textbf{GAP-ZC88-01} (which layer owns the step?).
  \item \textbf{ZC89} \textcolor{crfgreen}{$\checkmark$ CLOSES GAP-ZC88-01}
    via layer-specificity (THM-ZC89-01).
  \item \textbf{ZC90} opens \textbf{GAP-ZC90-01} (are the two mind-steps
    distinct kinds?) with CONJ-ZC90-01 attached.
  \item \textbf{ZC91} \textcolor{crfgreen}{$\checkmark$ CLOSES GAP-ZC90-01}
    and $\star$ RETIRES CONJ-ZC90-01 (irrationality-of-kind, THM-ZC91-01);
    opens \textbf{GAP-ZC91-01} (how do forward choices aggregate?).
  \item \textbf{ZC92} \textcolor{crfgreen}{$\approx$ PARTIALLY CLOSES
    GAP-ZC91-01} (Barami functional, common rate); opens \textbf{GAP-ZC92-01}
    (what gates the step kind?).
  \item \textbf{ZC93} \textcolor{crfgreen}{$\approx$ PARTIALLY CLOSES
    GAP-ZC92-01} (gate-ratio, THM-ZC93-01); the residual is handed forward
    beyond Part L.
\end{itemize}

\medskip
\textbf{Open at the close of Part L:} GAP-ZC91-01 (residual after ZC92's
partial close) and GAP-ZC92-01 (residual after ZC93's partial close). These
are the designated starting points for the next paper.
\end{correctionbox}

\end{document}
